
%
%
%
%
%
%
\chapter{Natural Semantics and Structural Operational Semantics: Formalization and Equivalence}
\label{formal_ns_sos_equiv}
First steps before reasoning whether Natural Semantics formalization and the Structural Operational Semantics formalization are equivalent is to create them.
In order to create them we also need a language to apply them to, in our case this will be \While \citep[p.~7]{nielson_semantics_2007}.


%
%
%
\section{Formalizing \While}
\label{formal_ns_sos_equiv:sec:while:formalization}

The first step in formalizing the equivalence between semantics is to
mechanize the language itself.
Due to Lean 4 \citep{theorem_proving_2026} being a programming language,
terms must be defined before they are used, which occasionally forces a
different order than traditional textbook presentations.

\subsection*{Syntactic Categories}

We begin by mechanizing the syntactic categories defined in chapter 1.2 of the textbook~\citep[p.~7]{nielson_semantics_2007}. 
The categories for \While are defined using BNF notation, mapping numerals ($n$), variables ($x$), and expressions ($a, b$) to inductive sets.
Using Lean 4's \texttt{inductive} keyword, we translate each BNF rule directly into a constructor:
\begin{paracol}{2}
    \begin{leftcolumn}
        \centering \textbf{Mechanized (Lean 4)}
    \end{leftcolumn}
    \begin{rightcolumn}
        \centering \textbf{Pen-and-Paper (Textbook)} \\\noindent
    \end{rightcolumn}
\end{paracol}
\vspace{1em}

\begin{paracol}{2}

    \switchcolumn[1]*[]
    \begin{leftcolumn}
        \begin{Lean_Code}
def Var := String
deriving instance DecidableEq for Var
        \end{Lean_Code}
    \end{leftcolumn}
    \begin{rightcolumn}
We assume that the structure of variables is given elsewhere;
for example, variables might be strings
of letters and digits starting with a letter.
    \end{rightcolumn}
\end{paracol}
\vspace{1em}

\begin{paracol}{2}

\switchcolumn[1]*[]
\begin{leftcolumn}
    \begin{Lean_Code}
inductive Num where
  | zero : Num
  | one  : Num
  | succ0 (n : Num) : Num
  | succ1 (n : Num) : Num
    \end{Lean_Code}
\end{leftcolumn}
\begin{rightcolumn}
    \(\displaystyle
\begin{array}{lcl}
    \\[0.3em]
    n &::=& 0 \\ 
        &\mid& 1 \\ 
        &\mid& n\ 1 \\ 
        &\mid& n\ 0
\end{array}
    \)
\end{rightcolumn}
\end{paracol}
\vspace{1em}

\begin{paracol}{2}
\begin{leftcolumn}
    \begin{Lean_Code}
inductive Aexp where
  | num (n : Num)
  | var (x : Var)
  | add (a1 a2 : Aexp)
  | mul (a1 a2 : Aexp)
  | sub (a1 a2 : Aexp)

-- Addition: @\color{green!40!black}\(a_1 + a_2\)@
-- (fullwidth plus \plus)
infixl:65 " @\color{purple}\(+\)@ " =>
    While.Aexp.add

-- Multiplication: @\color{green!40!black}\(a_1 \star a_2\)@
-- (star \star)
infixl:70 " @\color{purple}\(\star\)@ " =>
    While.Aexp.mul

-- Subtraction: @\color{green!40!black}\(a_1 - a_2\)@
-- (fullwidth minus \minus)
infixl:65 " @\color{purple}\(-\)@ " =>
    While.Aexp.sub
    \end{Lean_Code}
\end{leftcolumn}
\begin{rightcolumn}
    \(\displaystyle
\begin{array}{lcl}
    \\[0.3em]
    a &::=& n \\ 
        &\mid& x \\ 
        &\mid& a_1 + a_2 \\ 
        &\mid& a_1 \star a_2 \\ 
        &\mid& a_1 - a_2 \\
\end{array}
    \)
\end{rightcolumn}
\end{paracol}
\vspace{1em}

\begin{paracol}{2}
\begin{leftcolumn}
  \begin{Lean_Code}
inductive Bexp where
  | @true@
  | @false@
  | eq  (a1 a2 : Aexp)
  | le  (a1 a2 : Aexp)
  | not (b1 : Bexp)
  | and (b1 b2 : Bexp)

-- Equality: @\color{green!40!black}\(a_1\)@ == @\color{green!40!black}\(a_2\)@
-- (double equals ==)
notation:50
    a@\(_1\)@:51 " == " a@\(_2\)@:50 =>
    While.Bexp.eq @\(a_1\ a_2\)@

-- Less or equal: a@\color{green!40!black}\(\_1\)@ @\color{green!40!black}\(\leq\)@ a@\color{green!40!black}\(\_2\)@ 
-- (less or equal \leq)
notation:50
    a@\(_1\)@:51 " @\color{purple}\(\leq\)@ " a@\(_2\)@:50 =>
    While.Bexp.le @\(a_1\ a_2\)@

-- Not: @\color{green!40!black}\(\neg\)@b
prefix:75 "@\color{purple}\(\neg\)@" =>
    While.Bexp.not

-- And: b@\(_1\)@ @\color{green!40!black}\(\land\)@ b@\(_2\)@
infixr:35 " @\color{purple}\(\land\)@ " =>
    While.Bexp.and 

-- true/false
notation "@\color{purple}\(\mathbb{T}\)@" =>
    While.Bexp.true
notation "@\color{purple}\(\mathbb{F}\)@" =>
    While.Bexp.false
    \end{Lean_Code}
\end{leftcolumn}
\begin{rightcolumn}
    \(\displaystyle
\begin{array}{lcl}
    \\[0.3em]
    b &::=& \mathbf{true} \\ 
        &\mid& \mathbf{false} \\ 
        &\mid& a_1 = a_2 \\ 
        &\mid& a_1 \leq a_2 \\ 
        &\mid& \neg b_1 \\ 
        &\mid& b_1 \land b_2
\end{array}
    \)
\end{rightcolumn}
\end{paracol}
\vspace{1em}

\begin{paracol}{2}
\begin{leftcolumn}
    \begin{Lean_Code}
inductive Stmt where
  | ass  (x : Var)  (a : Aexp)
  | skip
  | composition (S@$_{1}$@ S@$_{2}$@ : Stmt)
  | cond (b : Bexp)
    (S@$_{1}$@ S@$_{2}$@ : Stmt)
  | loop (b : Bexp) (S@$_{1}$@ : Stmt)

-- Composition: S@\color{green!40!black}\(_1\)@; S@\color{green!40!black}\(_2\)@
infixl:80 "; " =>
    While.Stmt.composition

-- Assignment: x :@\color{green!40!black}\(\equiv\)@ a
-- (colon equiv)
infixr:90 " :@\color{purple}\(\equiv\)@ " =>
    While.Stmt.ass

-- Condition:
-- if b then S@\color{green!40!black}\(_1\)@ else S@\color{green!40!black}\(_2\)@
notation
    "if " b:40 " then "
    S@\(_1\)@:40 " else " S@\(_2\)@:40 =>
    While.Stmt.cond b S@\(_1\)@ S@\(_2\)@

-- While loop:
-- while b then S
notation
    "while " b:40 " then "
    S:40 =>
    While.Stmt.loop b S
    \end{Lean_Code}
\end{leftcolumn}
\begin{rightcolumn}
    \(\displaystyle
\begin{array}{lcl}
    \\[0.3em]
    S &::=& x := a \\
      &\mid& \mathtt{skip} \\
      &\mid& S_1; S_2 \\
      &\mid& \mathtt{if}\ b\ \mathtt{then}\ S_1\ \mathtt{else}\ S_2 \\
      \\
      &\mid& \mathtt{while}\ b\ \mathtt{do}\ S
\end{array}
    \)
\end{rightcolumn}
\end{paracol}
\vspace{1em}

\subsubsection*{Comparison of Syntactic Categories}
The translation of BNF to Lean's inductive types is largely direct, with some key differences:
\begin{itemize}
    \item \textbf{Implicit Definitions and Type Class Instances}:
    The textbook provides an informal description of variables ($x \in \mathbf{Var}$)
    as strings of letters and digits. In Lean, this must be a concrete \texttt{structure}.
    Crucially, to satisfy the Lean kernel's need for decidability (e.g., to
    determine if two variables in an assignment are the same), the type must
    explicitly derive \texttt{DecidableEq} and \texttt{BEq}, properties that
    are assumed as "given" in pen-and-paper proofs.

    \item \textbf{Structural Reformulation of Numerals}:
    While the textbook presents numerals ($n$) using a BNF that mimics string
    concatenation (e.g., $n\,0$, $n\,1$), the mechanized version uses a structured
    \texttt{inductive} type (\texttt{zero}, \texttt{one}, \texttt{succ0}, \texttt{succ1}).
    This is a significant structural shift: the textbook treats numerals as syntactic
    strings, while Lean treats them as a tree-structured data type to facilitate
    structural induction.

    \item \textbf{Explicit Operator Precedence}:
    The textbook introduces notation like \\
    "$\mathtt{if}\ b\ \mathtt{then}\ S_1\ \mathtt{else}\ S_2;\ S_3$" and relies on the
    reader's intuition or "standard" convention for disambiguation.
    Lean requires explicit binding priorities (e.g., \texttt{infixl:80}) to resolve
    whether the composition binds tighter than the conditional. This forces the
    formalizer to make "pedagogical shortcuts" regarding precedence explicit.

    \item \textbf{Explicit Typing and Domain Separation}:
    Each constructor in Lean explicitly specifies its input types
    (e.g., \texttt{Aexp} vs \texttt{Bexp}). In informal BNF, the membership of a
    term in a category is often inferred from context or the metavariable used
    ($a, b, S$), whereas the Lean kernel enforces these boundaries strictly
    during term construction.

    \item \textbf{Keyword and Namespace Constraints}:
    Practical mechanization requires navigating the host language's reserved
    keywords. We rename \texttt{if} to \texttt{cond} and \texttt{while} to
    \texttt{loop} within the \texttt{Stmt} inductive type to avoid collisions
    with Lean's internal control flow, a constraint entirely absent in the
    mathematical meta-language.
\end{itemize}
\noindent
With the syntactic structure of \While now formalized, we turn our attention to the semantic components,
the domains and evaluation functions that give meaning to these syntactic terms.

\subsection*{Semantic Domains and Functions}

With syntax formalized, we now specify the semantic domains and evaluation functions that interpret programs. These include the state representation and the functions that evaluate numerals, arithmetic expressions, and boolean expressions.

The \textbf{State} represents memory as a total function from variables to integers~\citep[p.~13]{nielson_semantics_2007}:
\begin{paracol}{2}
% --- SECTION: SEMANTIC DOMAINS ---
\begin{leftcolumn}
    \centering \textbf{Mechanized (Lean 4)}
\end{leftcolumn}
\begin{rightcolumn}
    \centering \textbf{Pen-and-Paper (Textbook)} \\\noindent
\end{rightcolumn}

\switchcolumn[1]*[]
\begin{leftcolumn}
    \begin{Lean_Code}
def State := Var @\(\to\)@ @ \(\mathbb{Z}\)@
    \end{Lean_Code}
\end{leftcolumn}
\begin{rightcolumn}
\(\displaystyle
s : \mathbf{Var} \to \mathbb{Z}
\)
\end{rightcolumn}

\switchcolumn[1]*[]
\begin{leftcolumn}
    \begin{Lean_Code}
def assign (s : State)
  (x : Var) (z : @\(\mathbb{Z}\)@) : State
    :=
  fun (v: Var) => 
    if v = x then z else s v

notation
  s "[" x " @\color{purple}\mapsto@ " z "]" => 
  assign s x z
    \end{Lean_Code}
\end{leftcolumn}
\begin{rightcolumn}
    \(\displaystyle
    s[x \mapsto v]y = \begin{cases} 
    v & \text{if } y = x \\ 
    s\,y & \text{if } y \neq x 
    \end{cases}
    \)
\end{rightcolumn}
\end{paracol}
\vspace{1em}

We now define the semantic evaluation functions $\mathcal{N}$, $\mathcal{A}$, and $\mathcal{B}$,
which map syntactic categories to their mathematical interpretations \missingcitation \todo{Cite the source for these specific denotations/naming conventions}:
\begin{paracol}{2}
% --- SECTION: SEMANTIC FUNCTIONS ---
\begin{leftcolumn}
    \centering \textbf{Mechanized (Lean 4)}
\end{leftcolumn}
\begin{rightcolumn}
    \centering \textbf{Pen-and-Paper (Textbook)} \\\noindent
\end{rightcolumn}

\switchcolumn[1]*[]
\begin{leftcolumn}
    \begin{Lean_Code}
set_option quotPrecheck
  false in
set_option hygiene false in
notation "@\color{purple}\(\mathcal{N}[\![\)@" n "@\color{purple}\(]\!]\)@" =>
  Num_eval  n
    \end{Lean_Code}
\end{leftcolumn}
\begin{rightcolumn}
    \[
    \mathcal{N} : Num \to \mathbb{Z}
    \]
\end{rightcolumn}

\switchcolumn[1]*[]
\begin{leftcolumn}
    \begin{Lean_Code}
def Num_eval : Num @\(\to\)@ @\(\mathbb{Z}\)@
  | Num.zero    => 0
  | Num.one     => 1
  | Num.succ0 n => 2*@\(\mathcal{N}[\![\)@n@\(]\!]\)@
  | Num.succ1 n => 2*@\(\mathcal{N}[\![\)@n@\(]\!]\)@ + 1
    \end{Lean_Code}
\end{leftcolumn}
\begin{rightcolumn}
    \[
    \begin{array}{rcl}
        \\[0.3em]
        \mathcal{N}[\![0]\!]s &=& 0 \\
        \mathcal{N}[\![1]\!]s &=& 1 \\
        \mathcal{N}[\![n\ 0]\!]s &=& 2 \cdot \mathcal{N}[\![n]\!] \\
        \mathcal{N}[\![n\ 1]\!]s &=& 2 \cdot \mathcal{N}[\![n]\!] + 1
    \end{array}
    \]
\end{rightcolumn}
\end{paracol}
\vspace{1em}

\begin{paracol}{2}
\begin{leftcolumn}
    \begin{Lean_Code}
set_option quotPrecheck
  false in
set_option hygiene false in
notation "@\color{purple}\(\mathcal{A}[\![\)@" a "@\color{purple}\(]\!]\)@" s:max =>
  Aexp_eval a s
    \end{Lean_Code}
\end{leftcolumn}
\begin{rightcolumn}
    \[
    \mathcal{A} : Aexp \to (State \to \mathbb{Z})
    \]
\end{rightcolumn}

\switchcolumn[1]*[]
\begin{leftcolumn}
    \begin{Lean_Code}
def Aexp_eval : Aexp @\(\to\)@
  (State @\(\to\)@ @\(\mathbb{Z}\)@)
  | Aexp.num n, _ => @\(\mathcal{N}[\![\)@n@\(]\!]\)@
  | Aexp.var x, s => s x
  | Aexp.add @\(a_1\)@ @\(a_2\)@, s =>
    @\(\mathcal{A}[\![a_1]\!]\)@s + @\(\mathcal{A}[\![a_2]\!]\)@s
  | Aexp.mul @\(a_1\)@ @\(a_2\)@, s =>
    @\(\mathcal{A}[\![a_1]\!]\)@s * @\(\mathcal{A}[\![a_2]\!]\)@s
  | Aexp.sub @\(a_1\)@ @\(a_2\)@, s =>
    @\(\mathcal{A}[\![a_1]\!]\)@s - @\(\mathcal{A}[\![a_2]\!]\)@s
    \end{Lean_Code}
\end{leftcolumn}
\begin{rightcolumn}
    \[
    \begin{array}{rcl}
        \\[1.8em]
        \mathcal{A}[\![n]\!]s &=& \mathcal{N}[\![n]\!] \\
        \mathcal{A}[\![x]\!]s &=& s \, x \\
        \mathcal{A}[\![a_1 + a_2]\!]s &=& \mathcal{A}[\![a_1]\!]s + \mathcal{A}[\![a_2]\!]s \\[1.15em]
        \mathcal{A}[\![a_1 \star a_2]\!]s &=& \mathcal{A}[\![a_1]\!]s \cdot \mathcal{A}[\![a_2]\!]s \\[1.15em]
        \mathcal{A}[\![a_1 - a_2]\!]s &=& \mathcal{A}[\![a_1]\!]s - \mathcal{A}[\![a_2]\!]s
    \end{array}
    \]
\end{rightcolumn}
\end{paracol}
\vspace{1em}

\begin{paracol}{2}
\begin{leftcolumn}
    \begin{Lean_Code}
set_option quotPrecheck
  false in
set_option hygiene false in
notation "@\color{purple}\(\mathcal{B}[\![\)@" b "@\color{purple}\(]\!]\)@" s:max =>
  Bexp_eval b s
    \end{Lean_Code}
\end{leftcolumn}
\begin{rightcolumn}
    \[
    \mathcal{B} : Bexp \to (State \to Bool)
    \]
\end{rightcolumn}

\switchcolumn[1]*[]
\begin{leftcolumn}
    \begin{Lean_Code}
def Bexp_eval : Bexp @\(\to\)@
  (State @\(\to\)@ Bool)
  | Bexp.@true@, _ => true
  | Bexp.@false@, _ => false
    \end{Lean_Code}
\end{leftcolumn}
\begin{rightcolumn}
    \[
    \begin{array}{rcl}
    \\[1.8em]
    \mathcal{B}[\![true]\!]s &=& \textbf{tt} \\
    \mathcal{B}[\![false]\!]s &=& \textbf{ff}
    \end{array}
    \]
\end{rightcolumn}

\switchcolumn[1]*[]
\begin{leftcolumn}
    \begin{Lean_Code}
  | Bexp.eq @\(a_1\)@ @\(a_2\)@, s =>
    @\(\mathcal{A}[\![a_1]\!]\)@s == @\(\mathcal{A}[\![a_2]\!]\)@s
    \end{Lean_Code}
\end{leftcolumn}
\begin{rightcolumn}
    \[
    \begin{array}{rcl}
    \mathcal{B}[\![a_1 = a_2]\!]s &=& \begin{cases} 
        \textbf{tt} & \text{if } \mathcal{A}[\![a_1]\!]s = \mathcal{A}[\![a_2]\!]s \\ 
        \textbf{ff} & \text{if } \mathcal{A}[\![a_1]\!]s \neq \mathcal{A}[\![a_2]\!]s
        \end{cases}
    \end{array}
    \]
\end{rightcolumn}

\switchcolumn[1]*[]
\begin{leftcolumn}
    \begin{Lean_Code}
  | Bexp.le @\(a_1\)@ @\(a_2\)@, s =>
    @\(\mathcal{A}[\![a_1]\!]\)@s @\(\leq\)@ @\(\mathcal{A}[\![a_2]\!]\)@s
    \end{Lean_Code}
\end{leftcolumn}
\begin{rightcolumn}
    \[
    \begin{array}{rcl}
    \mathcal{B}[\![a_1 \leq a_2]\!]s &=& \begin{cases} 
        \textbf{tt} & \text{if } \mathcal{A}[\![a_1]\!]s \leq \mathcal{A}[\![a_2]\!]s  \\ 
        \textbf{ff} & \text{if } \mathcal{A}[\![a_1]\!]s > \mathcal{A}[\![a_2]\!]s 
        \end{cases}
    \end{array}
    \]
\end{rightcolumn}

\switchcolumn[1]*[]
\begin{leftcolumn}
    \begin{Lean_Code}
  | Bexp.not @\(b\)@, s =>
    @\(\neg \mathcal{B}[\![b]\!]\)@s
    \end{Lean_Code}
\end{leftcolumn}
\begin{rightcolumn}
    \[
    \begin{array}{rcl}
    \mathcal{B}[\![\neg b]\!]s &=& \begin{cases} 
        \textbf{tt} & \text{if } \neg \mathcal{B}[\![b]\!]s \\ 
        \textbf{ff} & \text{if } \mathcal{B}[\![b]\!]s
        \end{cases}
    \end{array}
    \]
\end{rightcolumn}

\switchcolumn[1]*[]
\begin{leftcolumn}
    \begin{Lean_Code}
  | Bexp.and @\(b_1\)@ @\(b_2\)@, s =>
    @\(\mathcal{B}[\![b_1]\!]\)@s @\(\land\)@ @\(\mathcal{B}[\![b_2]\!]\)@s
    \end{Lean_Code}
\end{leftcolumn}
\begin{rightcolumn}
    \[
    \begin{array}{rcl}
    \mathcal{B}[\![b_1 \land b_2]\!]s &=& \begin{cases} 
        \textbf{tt} & \text{if } \mathcal{B}[\![b_1]\!]s = \textbf{tt} \\ & \text{and } \mathcal{B}[\![b_2]\!]s = \textbf{tt} \\ 
        \textbf{ff} & \text{if } \mathcal{B}[\![b_1]\!]s = \textbf{ff} \\ & \text{or } \mathcal{B}[\![b_2]\!]s = \textbf{ff} \\ 
        \end{cases}
    \end{array}
    \]
\end{rightcolumn}
\end{paracol}
\vspace{1em}


\subsubsection*{Comparison of Semantic Functions}
Formalizing evaluation functions $\mathcal{N}$, $\mathcal{A}$, and $\mathcal{B}$ reveals key differences between mechanized and pen-and-paper approaches:
\begin{itemize}
    \item \textbf{Explicit Notation}:
    While the textbook can directly introduce the notation of semantic brackets,
    the mechanized version must define this explicitly. To use notation within
    the definition itself (recursive calls), \texttt{quotPrecheck} and \texttt{hygiene}
    must be disabled to bypass Lean's requirement that identifiers be defined
    prior to their use in notation.

    \item \textbf{Boolean Realism}:
    Using built-in \texttt{Bool} (\texttt{true}/\texttt{false}) rather than
    symbolic constants $\mathbf{tt}$ and $\mathbf{ff}$ enables the use of
    Lean's internal boolean logic theorems and decision procedures,
    significantly reducing manual proof effort compared to symbolic manipulation.

    \item \textbf{Computational Case Analysis}:
    Where the textbook uses meta-language "if-then-else" blocks to define
    semantic values (e.g., for equality or less-than), Lean uses computational
    operators (like \texttt{==} or \texttt{$\leq$}). This collapses the
    distinction between the \textit{definition} of the operator and its \textit{evaluation}.

    \item \textbf{Explicit Totality}:
    Lean enforces that every function be total by requiring coverage of
    all constructors.

    \item \textbf{Type Consistency}:
    Lean's type system forces strict return types (e.g., $\mathbb{Z}$
    for $\mathcal{A}$ and \texttt{Bool} for $\mathcal{B}$), whereas
    pen-and-paper notation sometimes glosses over the boundary between
    the semantic domain and the meta-logic used to describe it.
\end{itemize}
\noindent
Beyond these core evaluation functions, mechanized reasoning about
programs requires formal properties of state manipulation that
pen-and-paper proofs often treat implicitly.

\subsection*{Auxiliary Proofs}

Beyond evaluation functions, state operations require formal verification. In addition to the above definitions, we need proofs to reason about state substitution:
First, we prove that repeated assignments to the same variable can be collapsed:
\[
    s[x \mapsto z_1][x \mapsto z_2] = s[x \mapsto z_2]
\]
This proof allows us to remove redundant assignments during semantic reasoning \missingcitation \todo{Cite a source for these state transformation lemmas, often found in 'Semantics with Applications'}.
The informal proof is straightforward by extensionality: two states are equal if they agree on all variables.
By definition of state assignment, both $s[x \mapsto z_1][x \mapsto z_2]$ and $s[x \mapsto z_2]$
map $x$ to $z_2$ and any other variable $y \neq x$ to $s(y)$.

\noindent
In Lean, this is formalized as:
\begin{Lean_Code}
theorem assign_shadow (s : State) (x : Var) (z1 z2 : @\(\mathbb{Z}\)@) :
  s[x @\mapsto@ z1][x @\mapsto@ z2] = s[x @\mapsto@ z2] := by @\dots@
\end{Lean_Code}

Second, we prove the commutativity of assignments to distinct variables:
\[
    s[x \mapsto z_1][y \mapsto z_2] = s[y \mapsto z_2][x \mapsto z_1], \quad y \neq x
\]

This allows us to reorder assignments to different variables, further simplifying state manipulation proofs.
The proof is again by extensionality: when $x \neq y$, both sides assign $z_1$ to $x$, $z_2$ to $y$,
and leave all other variables unchanged.
\begin{Lean_Code}
theorem assign_comm (s : State) (x y : Var) 
                       (z1 z2 : @\(\mathbb{Z}\)@) (h : x @\neq@ y) :
  s[x @\mapsto@ z1][y @\mapsto@ z2] = s[y @\mapsto@ z2][x @\mapsto@ z1] := by @\dots@
\end{Lean_Code}

Full proofs for these auxiliary theorems can be found in in the Appendix.
With these properties in place, apparently redundant or reordered assignments can be formally simplified, e.g.:
\[
    s[x \mapsto 3][x \mapsto 2][x \mapsto 1][x \mapsto 0] = s[x \mapsto 0]
\]

Having established the syntactic categories, semantic domains, and evaluation functions,
we now possess the foundation needed to define the transition systems.
The next section formalizes Natural Semantics (also called big-step semantics) for the \While language,
followed by Structural Operational Semantics (small-step), and finally establishes their equivalence.

%
%
%
\section{Formalizing Natural Semantics (Big-Step)}
\label{formal_ns_sos_equiv:sec:ns:formalization}
Now that we have a mechanized definition of the \While language we can proceed by defining the actual semantics of running a \While language program. 
We will start by defining the natural operational semantics often referred to as "big step semantics" \missingcitation \todo{Cite Gilles Kahn (1987) or Nielson \& Nielson) for the introduction of Natural Semantics}.
When it comes to operational semantics we use deduction tree syntax to provide rules and axioms.

\begin{paracol}{2}
\begin{leftcolumn}
    \centering \textbf{Mechanized (Lean 4)}
\end{leftcolumn}
\begin{rightcolumn}
    \centering \textbf{Pen-and-Paper (Textbook)} \\\noindent
\end{rightcolumn}

\switchcolumn[1]*[]
\begin{leftcolumn}
    \begin{Lean_Code}
set_option quotPrecheck
  false in
set_option hygiene false in
notation:40
  "@\color{purple}\(\langle\)@" S "," s "@\color{purple}\(\rangle\)@" " @\color{purple}\(\to_{ns}\)@ " s'
  => big_step S s s'

inductive big_step : Stmt @\(\to\)@
  State @\(\to\)@ State @\(\to\)@ Prop where
    \end{Lean_Code}
\end{leftcolumn}
\begin{rightcolumn}
\end{rightcolumn}
\vspace{2em}
\switchcolumn[1]*[]
\begin{leftcolumn}
    \begin{Lean_Code}
  | ass {s x a} :
      @\(\langle\)@x @\(: \equiv\)@ a, s@\(\rangle\)@ @\(\to_{ns}\)@
      (s[x @\mapsto@ @\(\mathcal{A}[\![a]\!]\)@ s])
\end{Lean_Code}
\end{leftcolumn}
\begin{rightcolumn}
\vspace{1.1em}
    \(
\begin{prooftree}
  \hypo{} % Empty hypothesis for an axiom
  \infer[
    left label=[ass$_{\text{ns}}$]
    ]1{
        \parbox{2.5cm}{$
    \langle x := a, s \rangle \to \\
    s[x \mapsto \mathcal{A}[\![a]\!] s]
        $}
    }
\end{prooftree}
            \)
\end{rightcolumn}
\vspace{1em}
\switchcolumn[1]*[]
\begin{leftcolumn}
    \begin{Lean_Code}
  | skip {s} :
      @\(\langle\)@Stmt.skip, s@\(\rangle\)@ @\(\to_{ns}\)@ s
        \end{Lean_Code}
\end{leftcolumn}
\begin{rightcolumn}
\vspace{1.1em}
    \(
\begin{prooftree}
  \hypo{} % Empty hypothesis for an axiom
  \infer[
    left label=[skip$_{\text{ns}}$]
    ]1{
        \langle \mathtt{skip}, s \rangle \to s
    }
\end{prooftree}
    \)
\end{rightcolumn}
\end{paracol}
\newpage

\begin{paracol}{2}
\begin{leftcolumn}
    \begin{Lean_Code}
  | comp {S@$_{1}$@ S@$_{2}$@ s s' s''}
    (h_left  : @\(\langle\)@S@\(_1\)@, s@\(\rangle\)@ @\(\to_{ns}\)@ s')
    (h_right :
      @\(\langle\)@S@\(_2\)@, s'@\(\rangle\)@ @\(\to_{ns}\)@ s'') :
      @\(\langle\)@S@\(_1\)@; S@\(_2\)@, s@\(\rangle\)@ @\(\to_{ns}\)@ s''
    \end{Lean_Code}
\end{leftcolumn}
\begin{rightcolumn}
\vspace{3.6em}
    \(
\begin{prooftree}
  \hypo{\langle S_1, s \rangle \to s'}
  \hypo{\langle S_2, s' \rangle \to s''}
  \infer[
    left label=[comp$_{\text{ns}}$]
    ]2{
        \langle S_1;\ S_2, s \rangle \to s''
    }
\end{prooftree}
    \)
\end{rightcolumn}
\vspace{3em}
\switchcolumn[1]*[]
\begin{leftcolumn}
    \begin{Lean_Code}
  | if_true {b S@$_{1}$@ S@$_{2}$@ s s'}
    (h_cond : @\(\mathcal{B}[\![b]\!]\)@ s = true)
    (h_then : @\(\langle\)@S@\(_1\)@, s@\(\rangle\)@ @\(\to_{ns}\)@ s') :
    @\(\langle\)@if b then S@$_{1}$@ else S@$_{2}$@, s@\(\rangle\)@
    @\(\to_{ns}\)@ s'
    \end{Lean_Code}
\end{leftcolumn}
\begin{rightcolumn}
\vspace{2.5em}
    \(
\begin{prooftree}
  \hypo{\langle S_1, s \rangle \to s'}
  \infer[
    left label={[if$_{\text{ns}}^{\text{\ tt}}$]}
    ]1[
    if $\mathcal{B}[\![b]\!]s = \mathbf{tt}$
    ]{
        \parbox{4cm}{$
    \langle \mathtt{if}\ b\ \mathtt{then}\ S_1\ \mathtt{else}\ S_2, s \rangle \\
    \to s'
        $}
    }
\end{prooftree}
    \)
\end{rightcolumn}
\vspace{3em}
\switchcolumn[1]*[]
\begin{leftcolumn}
    \begin{Lean_Code}
  | if_false {b S@$_{1}$@ S@$_{2}$@ s s'}
    (h_cond : @\(\mathcal{B}[\![b]\!]\)@ s = false)
    (h_else : @\(\langle\)@S@\(_2\)@, s@\(\rangle\)@ @\(\to_{ns}\)@ s') :
    @\(\langle\)@if b then S@$_{1}$@ else S@$_{2}$@, s@\(\rangle\)@
    @\(\to_{ns}\)@ s'
    \end{Lean_Code}
\end{leftcolumn}
\begin{rightcolumn}
\vspace{2.5em}
    \(
\begin{prooftree}
  \hypo{\langle S_1, s \rangle \to s'}
  \infer[
    left label={[if$_{\text{ns}}^{\text{\ ff}}$]}
    ]1[
    if $\mathcal{B}[\![b]\!]s = \mathbf{ff}$
    ]{
        \parbox{4cm}{$
    \langle \mathtt{if}\ b\ \mathtt{then}\ S_1\ \mathtt{else}\ S_2, s \rangle \\
    \to s'
        $}
    }
\end{prooftree}
    \)
\end{rightcolumn}
\vspace{3em}
\switchcolumn[1]*[]
\begin{leftcolumn}
    \begin{Lean_Code}
  | while_true {b S s s' s''}
    (h_cond : @\(\mathcal{B}[\![b]\!]\)@ s = true)
    (h_step : @\(\langle\)@S, s@\(\rangle\)@ @\(\to_{ns}\)@ s')
    (h_rest :
      @\(\langle\)@while b then S, s'@\(\rangle\)@
      @\(\to_{ns}\)@ s'') :
      @\(\langle\)@while b then S, s@\(\rangle\)@
      @\(\to_{ns}\)@ s''
    \end{Lean_Code}
\end{leftcolumn}
\begin{rightcolumn}
\vspace{5.4em}
    \(
\begin{prooftree}
  \hypo{\langle S, s \rangle \to s'}
  \hypo{
        \parbox{3cm}{$
    \langle \mathtt{while}\ b\ \mathtt{do}\ S, s' \rangle \\
    \to s''
        $}
    }
  \infer[
      left label={[while$_{\text{ns}}^{\text{\ tt}}$]}
      ]2[
      if $\mathcal{B}[\![b]\!]s = \mathbf{tt}$
      ]{
        \parbox{5cm}{$
    \langle \mathtt{while}\ b\ \mathtt{do}\ S, s \rangle \\
    \to s''
        $}
    }
\end{prooftree}
    \)
\end{rightcolumn}
\vspace{2em}
\switchcolumn[1]*[]
\begin{leftcolumn}
    \begin{Lean_Code}
  | while_false {b S' s}
    (h_cond : @\(\mathcal{B}[\![b]\!]\)@ s = false) :
      @\(\langle\)@while b then S, s@\(\rangle\)@
      @\(\to_{ns}\)@ s
    \end{Lean_Code}
\end{leftcolumn}
\begin{rightcolumn}
    \vspace{2.5em}
    \(
\begin{prooftree}
  \hypo{} % Empty hypothesis for an axiom
  \infer[
    left label={[while$_{\text{ns}}^{\text{\ ff}}$]}
    ]1[
    if $\mathcal{B}[\![b]\!]s = \mathbf{ff}$
    ]{
        \parbox{3cm}{$
    \langle \mathtt{while}\ b\ \mathtt{do}\ S, s \rangle \\
    \to s
        $}
    }
\end{prooftree}
    \)
\end{rightcolumn}
\end{paracol}
\vspace{1em}

\subsubsection*{Comparison and Result}
The transition system for Natural Semantics is formalized as an inductive
relation \texttt{big\_step} that captures the complete execution of a \While
program from initial state to final state. This big-step approach mirrors the
pen-and-paper formulation: each rule represents either an axiom or a deductive rule.
However, the mechanization necessitates making several implicit textbook
conventions explicit.

\begin{itemize}
    \item \textbf{Direct Structural Mapping (RQ1)}:
    There is a near 1-to-1 mapping between the deduction tree rules in the
    curriculum and Lean's \texttt{inductive} constructors. The \texttt{big\_step}
    relation effectively treats the constructors as the inference rules themselves,
    confirming that the structural reformulation is minimal.
    
    \item \textbf{Explicit Side Conditions (RQ2)}:
    In pen-and-paper proofs, side conditions (e.g.,
    $if \ \mathcal{B}[\![b]\!]s = \mathbf{tt}$) are often treated as meta-level
    annotations to justify the choice of a rule. In Lean, these are
    "pedagogical shortcuts" that must be promoted to formal premises (\texttt{h\_cond}).
    The Lean kernel requires a proof of the boolean evaluation, whereas the curriculum
    assumes this evaluation happens "outside" the derivation tree.
    
    \item \textbf{Typing of Semantic Domains}:
    The textbook notation $\mathcal{B}[\![b]\!]s = \mathbf{tt}$ uses a boldface
    $\mathbf{tt}$ to distinguish the semantic value from the syntactic expression.
    In Lean, this requires an explicit mapping to the \texttt{Bool} type, forcing a
    distinction between the Prop-level equality and the boolean value returned by the
    semantic function.
    
    \item \textbf{Axiom vs. Rule Uniformity}:
    While the curriculum distinguishes between axioms (horizontal line with no hypothesis)
    and rules, Lean treats both uniformly as constructors. The "shortcut" in the literature
    is the visual omission of premises for \texttt{ass} and \texttt{skip},
    which in Lean simply translates to constructors that do not take other
    \texttt{big\_step} instances as arguments.
    
    \item \textbf{Metavariable Scope}:
    In the \texttt{comp} and \texttt{while\_true} rules, the intermediate state
    $s'$ is implicitly existentially quantified in the literature. Lean requires
    these intermediate states to be explicitly named and managed within the
    constructor's signature, making the data flow through the derivation tree
    entirely transparent.
\end{itemize}

\noindent
This mechanization preserves the meaning of the original rules while enforcing
that all premises must be discharged before the conclusion can be derived.
The notation $\langle S, s \rangle \to_{ns} s'$ is defined to match the
textbook's presentation, allowing proofs to read similarly to their
pen-and-paper counterparts.

With the transition system now formalized, we can prove properties about Natural Semantics.

%
%
%
\section{Formalizing Proofs in Natural Semantics}
\label{formal_ns_sos_equiv:sec:ns:proofs}
Proofs made for natural semantics in the literature tend to share a pattern in using induction over the shape of the tree.
More formally the book uses the following definition:
\begin{quotation}
\noindent
Induction on the Shape of Derivation Trees\vspace{1em} \\  \noindent
1: Prove that the property holds for all the simple derivation trees by
showing that it holds for the axioms of the transition system.\vspace{1em} \\ \noindent
2: Prove that the property holds for all composite derivation trees: For
each rule assume that the property holds for its premises (this is
called the induction hypothesis) and prove that it also holds for the
conclusion of the rule provided that the conditions of the rule are
satisfied.
\end{quotation}
Doing induction over the shape will be used in the proof in \autoref{formal_ns_sos_equiv:sec:ns:while_unroll}
Let's start with an easier proof without any induction.

%
%
\subsection{While-loop Unrolling (Lemma 2.5)}
\label{formal_ns_sos_equiv:sec:ns:while_unroll}
In the textbook~\citep[p.~26]{nielson_semantics_2007} they provide a lemma stating the following:
\begin{quote}
The statement
\[\mathtt{while}\ b\ \mathtt{do}\ S\]
is semantically equivalent to
\[\mathtt{if}\ b\ \mathtt{then} (S;\ \mathtt{while}\ b\ \mathtt{do}\ S)\ \mathtt{else}\ \mathtt{skip}\]
\end{quote}
This would be equivalent with stating the following:
\[
    \langle \mathtt{while}\ b\ \mathtt{do}\ S,\ s\rangle \to s'' \iff  \langle \mathtt{if}\ b\ \mathtt{then} (S;\ \mathtt{while}\ b\ \mathtt{do}\ S)\ \mathtt{else}\ \mathtt{skip},\ s\rangle \to s''
\]
Making sure to label them the same way they do in the textbook, we get the following:
\[
\langle \mathtt{while}\ b\ \mathtt{do}\ S,\ s\rangle \to s'' \tag{*}
\]
\[
\langle \mathtt{if}\ b\ \mathtt{then} (S;\ \mathtt{while}\ b\ \mathtt{do}\ S)\ \mathtt{else}\ \mathtt{skip},\ s\rangle \to s'' \tag{**}
\]
In the textbook they divide the proof into to two parts, first part is to prove that given (*) that (**) holds,
the second part proves the reverse, namely proving that given (**) that (*) holds. Which is the normal way to prove equivalence of two statements.

\subsubsection*{Forward Direction: $(*) \implies (**)$}
Let us now compare the textbook proof with the mechanized equivalent for the case where we show that (*) holds, implies that (**) holds.
\begin{paracol}{2}
\begin{leftcolumn}
    \centering \textbf{Mechanized (Lean 4)}
\end{leftcolumn}
\begin{rightcolumn}
    \centering \textbf{Pen-and-Paper (Textbook)} \\\noindent
\end{rightcolumn}

\switchcolumn[1]*[]
\begin{leftcolumn}
    \begin{Lean_Code}
theorem while_expand
  (b : Bexp)
  (S : Stmt)
  (s s' : State) :
  (@\(\langle\)@while b then S, s@\(\rangle\)@ @\(\to_{ns}\)@ s')
  @\(\to\)@ (@\(\langle\)@if b
    then (S; while b then S)
    else Stmt.skip, s@\(\rangle\)@ @\(\to_{ns}\)@ s')
  := by
    \end{Lean_Code}
\end{leftcolumn}
\begin{rightcolumn}
\end{rightcolumn}
\switchcolumn[1]*[]
\begin{leftcolumn}
    \begin{Lean_Code}
  -- Assume (*) holds
  -- from implication above
  intro h_while
    \end{Lean_Code}
\end{leftcolumn}
\begin{rightcolumn} \noindent
Because (*) holds, we know that we have a derivation tree T for it.
\end{rightcolumn}
\switchcolumn[1]*[]
\begin{leftcolumn}
    \begin{Lean_Code}
  -- Split into cases based on
  -- the last rule applied to
  -- derive (*)
  cases h_while with
    \end{Lean_Code}
\end{leftcolumn}
\begin{rightcolumn} \noindent
It can have one of two forms depending on whether it has been constructed using the
rule [while$_{\text{ns}}^{\text{\ tt}}$] or the axiom [while$_{\text{ns}}^{\text{\ ff}}$].
\end{rightcolumn}
\switchcolumn[1]*[]
\begin{leftcolumn}
    \begin{Lean_Code}
  -- Case: [while@\color{green!40!black}$_{\text{ns}}^{\text{\ tt}}$@]
  | while_true
    h_cond_true T@$_1$@ T@$_2$@ =>
    \end{Lean_Code}
\end{leftcolumn}
\begin{rightcolumn} \noindent
In the first case, the derivation tree T has the form
\[
    \begin{prooftree}
        \hypo{T_1}
        \hypo{T_2}
        \infer[
            ]2{
            \langle \mathtt{while}\ b\ \mathtt{do}\ S,\ s\rangle \to s''
            }
    \end{prooftree}
\]
where \(T_1\) is a derivation tree with root \(\langle S,\ s\rangle \to s'\) and \(T_2\) is a derivation tree
with root \(\langle \mathtt{while}\ b\ \mathtt{do}\ S,\ s\rangle \to s''\).
Furthermore, \(\mathcal{B}[\![b]\!]s = \mathbf{tt}\).
\end{rightcolumn}
\vspace{1.5em}

\switchcolumn[1]*[]
\begin{leftcolumn}
    \begin{Lean_Code}
    -- Construct the composition [comp@\color{green!40!black}$_{\text{ns}}$@]
    -- using T@\color{green!40!black}$_1$@ and T@\color{green!40!black}$_2$@
    have comp_stmt :
      (@\(\langle\)@S; while b then S, s@\(\rangle\)@
      @\(\to_{ns}\)@ s') :=
      big_step.comp T@$_1$@ T@$_2$@
    \end{Lean_Code}
\end{leftcolumn}

\begin{rightcolumn} \noindent
Using the derivation trees \(T_1\) and \(T_2\) as the premises for the rules [comp$_{\text{ns}}$], we can construct the derivation tree
\[
    \begin{prooftree}
        \hypo{T_1}
        \hypo{T_2}
        \infer[
            % left label=[comp$_{\text{ns}}$]
            ]2{
            \langle S;\ \mathtt{while}\ b\ \mathtt{do}\ S,\ s\rangle \to s''
            }
    \end{prooftree}
\]
\end{rightcolumn}
\vspace{1.5em}

\switchcolumn[1]*[]
\begin{leftcolumn}
    \begin{Lean_Code}
    -- Apply [if@\color{green!40!black}$_{\text{ns}}^{\ \text{tt}}$@] rule
    exact big_step.if_true
      h_cond_true
      comp_stmt
    \end{Lean_Code}
\end{leftcolumn}

\begin{rightcolumn} \noindent
Using that \(\mathcal{B}[\![b]\!]s = \mathbf{tt}\), we can use the rule [if$_{\text{ns}}^{\ \text{tt}}$]
to construct the derivation tree
\[
    \begin{prooftree}
        \hypo{T_1}
        \hypo{T_2}
        \infer[
            % left label=[comp$_{\text{ns}}$]
            ]2{
            \langle S;\ \mathtt{while}\ b\ \mathtt{do}\ S,\ s\rangle \to s''
            }
        \infer[
            % left label=[if$_{\text{ns}}^{\ \text{tt}}$]
            ]1{
            \langle \mathtt{if}\ b\ \mathtt{then}\ (S;\ \mathtt{while}\ b\ \mathtt{do}\ S)\ \mathtt{else}\ \mathtt{skip},\ s\rangle \to s''
            }
    \end{prooftree}
\]
hereby showing that (**) holds.
\end{rightcolumn}
\vspace{2em}
\switchcolumn[1]*[]
\begin{leftcolumn}
    \begin{Lean_Code}
  -- Case: [while@\color{green!40!black}$_{\text{ns}}^{\text{\ ff}}$@]
  | while_false h_cond_false =>
    \end{Lean_Code}
\end{leftcolumn}
\begin{rightcolumn} \noindent
Alternatively, the derivation tree \(T\) is an instance of [while$_{\text{ns}}^{\ \text{ff}}$].
Then \(\mathcal{B}[\![b]\!]s = \mathbf{ff}\) and we must have that \(s'' = s\). So \(T\) simply is
\[
    \begin{prooftree}
        \hypo{\langle \mathtt{while}\ b\ \mathtt{do}\ S,\ s\rangle \to s}
    \end{prooftree}
\]
\end{rightcolumn}
\vspace{1.5em}

\switchcolumn[1]*[]
\begin{leftcolumn}
    \begin{Lean_Code}
    -- Construct [skip@\color{green!40!black}$_{\text{ns}}$@] axiom
    have skip_stmt :
      (@\(\langle\)@Stmt.skip, s@\(\rangle\)@
      @\(\to_{ns}\)@ s) :=
      big_step.skip
    \end{Lean_Code}
\end{leftcolumn}
\begin{rightcolumn} \noindent
Using the axiom [skip$_{\text{ns}}$], we get a derivation tree
\[
    \begin{prooftree}
        \hypo{\langle \mathtt{skip},\ s\rangle \to s''}
    \end{prooftree}
\]
\end{rightcolumn}
\vspace{1.5em}

\switchcolumn[1]*[]
\begin{leftcolumn}
    \begin{Lean_Code}
    -- Apply [if@\color{green!40!black}$_{\text{ns}}^{\ \text{ff}}$@] rule
    exact big_step.if_false
      h_cond_false
      skip_stmt
    \end{Lean_Code}
\end{leftcolumn}

\switchcolumn
\begin{rightcolumn} \noindent
and we can now apply the rule [if$_{\text{ns}}^{\ \text{ff}}$] to construct a derivation tree for (**):
\[
    \begin{prooftree}
        \hypo{\langle \mathtt{skip},\ s\rangle \to s''}
        \infer[
        ]1{
            \langle \mathtt{if}\ b\ \mathtt{then}\ (S;\ \mathtt{while}\ b\ \mathtt{do}\ S)\ \mathtt{else}\ \mathtt{skip},\ s\rangle \to s''
        }
    \end{prooftree}
\]
This completes the first part of the proof.
\end{rightcolumn}
\end{paracol}
\vspace{1em}

As we can see the Lean versions mimics the textbook very closely.
This proof builds the deduction tree from the top and moving down towards the
root with our wanted result.


\subsubsection*{Backward Direction: $(**) \implies (*)$}
Now we will show that (**) holds, implies that (*) holds.
\begin{paracol}{2}
\begin{leftcolumn}
    \centering \textbf{Mechanized (Lean 4)}
\end{leftcolumn}
\begin{rightcolumn}
    \centering \textbf{Pen-and-Paper (Textbook)} \\\noindent
\end{rightcolumn}

\switchcolumn[1]*[]
\begin{leftcolumn}
    \begin{Lean_Code}
theorem if_compact (b : Bexp)
  (S : Stmt) (s s' : State) :
  (@\(\langle\)@if b
  then (S; while b then S)
  else Stmt.skip, s@\(\rangle\)@ @\(\to_{ns}\)@ s')
  @\(\to\)@
  (@\(\langle\)@while b then S, s@\(\rangle\)@ @\(\to_{ns}\)@ s')
  := by
    \end{Lean_Code}
\end{leftcolumn}
\begin{rightcolumn} \noindent
For the second part of the proof,
\end{rightcolumn}
\vspace{1.5em}

\switchcolumn[1]*[]
\begin{leftcolumn}
    \begin{Lean_Code}
  -- Assume (**) holds
  intro h_if
    \end{Lean_Code}
\end{leftcolumn}
\begin{rightcolumn} \noindent
we assume that (**) holds and shall prove
that (*) holds. So we have a derivation tree \(T\) for (**)
\end{rightcolumn}
\vspace{1.5em}

\switchcolumn[1]*[]
\begin{leftcolumn}
    \begin{Lean_Code}
  -- construct tree for (*)
  show (@\(\langle\)@while b then S, s@\(\rangle\)@
    @\(\to_{ns}\)@ s')
    \end{Lean_Code}
\end{leftcolumn}
\begin{rightcolumn} \noindent
and must construct one for (*).
\end{rightcolumn}
\vspace{1.5em}

\switchcolumn[1]*[]
\begin{leftcolumn}
    \begin{Lean_Code}
  -- Only [if@\color{green!40!black}$_{\text{ns}}^{\ \text{tt}}$@] or [if@\color{green!40!black}$_{\text{ns}}^{\ \text{ff}}$@]
  -- can have been used
  -- for (**)
  cases h_if with
    \end{Lean_Code}
\end{leftcolumn}
\begin{rightcolumn} \noindent
Only two rules could give rise to the derivation tree \(T\) for (**), namely
[if$_{\text{ns}}^{\text{\ tt}}$] or [if$_{\text{ns}}^{\text{\ ff}}$].
\end{rightcolumn}
\vspace{1.5em}

\switchcolumn[1]*[]
\begin{leftcolumn}
    \begin{Lean_Code}
  -- Case: [if@\color{green!40!black}$_{\text{ns}}^{\ \text{tt}}$@]
  | if_true
    h_cond_true T@$_1$@ =>
    \end{Lean_Code}
\end{leftcolumn}
\begin{rightcolumn} \noindent
In the first case, \(\mathcal{B}[\![b]\!]s = \mathbf{tt}\) and we have a derivation tree \(T_1\)
with root
\[
    \begin{prooftree}
        \hypo{\langle S;\ \mathtt{while}\ b\ \mathtt{do}\ S,\ s\rangle \to s''}
    \end{prooftree}
\]
\end{rightcolumn}
\vspace{1.5em}

\switchcolumn[1]*[]
\begin{leftcolumn}
    \begin{Lean_Code}
    -- Invert Composition
    -- [comp@\color{green!40!black}$_{\text{ns}}$@]
    cases T@$_1$@ with
    | comp T@$_2$@ T@$_3$@ =>
    \end{Lean_Code}
\end{leftcolumn}
\begin{rightcolumn} \noindent
The statement has the general form \(S_1;\ S_2\), and the only rule that could give
this is [comp$_{\text{ns}}$]. Therefore there are derivation trees \(T _2\) and \(T_3\) for
\[
    \begin{prooftree}
        \hypo{\langle S,\ s\rangle \to s'}
    \end{prooftree}
\]
and
\[
    \begin{prooftree}
        \hypo{\langle \mathtt{while}\ b\ \mathtt{do}\ S,\ s'\rangle \to s''}
    \end{prooftree}
\]
for some state \(s'\).
\end{rightcolumn}
\vspace{1.5em}

\switchcolumn[1]*[]
\begin{leftcolumn}
    \begin{Lean_Code}
      -- Combine into
      -- [while@\color{green!40!black}$_{\text{ns}}^{\text{\ tt}}$@]
      exact big_step.while_true
        h_cond_true
        T@$_2$@ T@$_3$@
    \end{Lean_Code}
\end{leftcolumn}
\begin{rightcolumn} \noindent
It is now straightforward to use the rule [while$_{\text{ns}}^{\text{\ tt}}$] to combine
\(T _2\) and \(T_3\) into a derivation tree for (*).
\end{rightcolumn}
\vspace{1.5em}

\switchcolumn[1]*[]
\begin{leftcolumn}
    \begin{Lean_Code}
  -- Case: [if@\color{green!40!black}$_{\text{ns}}^{\ \text{ff}}$@]
  | if_false
    h_cond_false
    h_deriv_skip =>
    \end{Lean_Code}
\end{leftcolumn}

\begin{rightcolumn} \noindent
In the second case, \(\mathcal{B}[\![b]\!]s = \mathbf{ff}\) and \(T\) is constructed using the rule [while$_{\text{ns}}^{\text{\ ff}}$].
\end{rightcolumn}
\vspace{1.5em}

\switchcolumn[1]*[]
\begin{leftcolumn}
    \begin{Lean_Code}
    -- Invert [skip@\color{green!40!black}$_{\text{ns}}$@] axiom
    cases h_deriv_skip
    \end{Lean_Code}
\end{leftcolumn}

\begin{rightcolumn} \noindent
This means that we have a derivation tree for
\[
    \begin{prooftree}
        \hypo{\langle \mathtt{skip},\ s\rangle \to s''}
    \end{prooftree}
\]
and according to axiom [skip$_{\text{ns}}$] it must be the case that \(s = s''\).
\end{rightcolumn}
\vspace{1.5em}

\switchcolumn[1]*[]
\begin{leftcolumn}
    \begin{Lean_Code}
    -- Apply [while@\color{green!40!black}$_{\text{ns}}^{\text{\ ff}}$@] rule
    exact big_step.while_false
      h_cond_false
    \end{Lean_Code}
\end{leftcolumn}

\begin{rightcolumn} \noindent
But then we can use the axiom [while$_{\text{ns}}^{\text{\ ff}}$] to construct a derivation tree for (*).
This completes the proof.
\end{rightcolumn}
\vspace{1.5em}

\end{paracol} 

\subsubsection*{Applying the Equivalence}
Whilst the textbook is now done with their proof, we must continue to formulate the equivalence relation,
this is done by using the previous two proofs in a new final proof constructed as following:
\begin{Lean_Code}
lemma while_unroll (b : Bexp) (S : Stmt) (s s' : State) :
  (@\(\langle\)@if b then (S; while b then S) else Stmt.skip, s@\(\rangle\)@ @\(\to_{ns}\)@ s')
  @\(\leftrightarrow\)@
  (@\(\langle\)@while b then S, s@\(\rangle\)@ @\(\to_{ns}\)@ s') := by
  -- equivalence introduction
  apply Iff.intro

  -- Forward direction: if => while
  case mp =>
    -- Reusing our previous proof for if => while
    exact if_compact b S s s'

  -- Backward direction: while => if
  case mpr =>
    -- Reusing our previous proof for while => if
    exact while_expand b S s s'
\end{Lean_Code}

\subsubsection*{Comparison and Result}
This proof demonstrates the modular nature of both formal and informal proofs.
By splitting the equivalence ($\iff$) into two implications
($\implies$ and $\impliedby$), we can prove each direction independently.
This modularity is mirrored in Lean through separate theorems
(while\_expand and if\_compact) which are then combined using Iff.intro.
Key observations include:

\begin{itemize}
    \item \textbf{Syntactic Correspondence}:
    The mechanized proof maps nearly 1:1 to the textbook's structure (RQ1).
    Each \texttt{case} in Lean corresponds to a branch in the manual proof,
    and each rule application (e.g., \texttt{big\_step.comp}) matches the
    "construction of a new tree" in the text.
    
    \item \textbf{Explicit vs. Implicit Construction}:
    In Lean, we use the \texttt{have} tactic to explicitly name intermediate
    derivation trees (like \texttt{comp\_stmt}). The textbook does this narratively
    ("we can construct the derivation tree...").
    
    \item \textbf{Exhaustive Matching}:
    Lean's \texttt{cases} tactic on the \texttt{if} statement automatically
    identifies that only \texttt{if\_true} and \texttt{if\_false} are possible,
    mirroring the textbook's assertion "Only two rules could give rise to...".
    
    \item \textbf{Implicit Rule Inversion (RQ2)}: A notable
    "pedagogical shortcut" in the textbook is the phrasing
    "it must be the case that $s = s''$" when examining the \texttt{skip} axiom.
    In Lean, this requires an explicit inversion or destruction of the
    \texttt{skip} hypothesis (\texttt{cases h\_deriv\_skip}) to unify the
    state variables in the goal.
    
    \item \textbf{State Unification}:
    The textbook often assumes variables like $s''$ are automatically
    "carried over" between trees. Lean requires explicit alignment of
    these state variables, making the implicit assumption that the state
    resulting from \texttt{skip} is identical to the starting state a
    formal necessity.
\end{itemize}


%
%
\subsection{Determinism (Theorem 2.9)}
\label{formal_ns_sos_equiv:sec:ns:determinism}
The next proof made in the textbook~\citep[p.~29]{nielson_semantics_2007} refers to determinism, where the definition of determinism is given
\begin{quotation}
\dots semantics is called deterministic if for all choices of
\(S\), \(s\), \(s'\) and \(s''\) we have that
\[
\langle S,\ s \rangle \to s'\ \mbox{and}\ \langle S,\ s \rangle \to s''\ \mbox{imply}\ s' = s''
\]
This means that for every statement \(S\) and initial state \(s\) we can uniquely
determine a final state \(s'\) if and only if the execution of \(S\) terminates.
\end{quotation}
By using this definition we can construct a proof using the following formulation:
\[
\langle S,\ s \rangle \to s' \implies \left(\langle S,\ s \rangle \to s'' \implies s' = s''\right)
\]
We will show the proofs side by side to show what the equivalent steps would be in Lean 4 compared to the textbook.
\vspace{1em}
\begin{paracol}{2}
    \begin{leftcolumn}
        \centering \textbf{Mechanized (Lean 4)}
    \end{leftcolumn}
    \begin{rightcolumn}
        \centering \textbf{Pen-and-Paper (Textbook)} \\\noindent
    \end{rightcolumn}
    \vspace{1.5em}


\switchcolumn[1]*[]
\begin{leftcolumn}
    \begin{Lean_Code}
theorem deterministic
  (S : Stmt)
  (s s' s'' : State)
  (h: @\(\langle\)@S, s@\(\rangle\)@ @\to$_{ns}$@ s') :
  (@\(\langle\)@ S, s@\(\rangle\)@ @\to$_{ns}$@ s'')
  @\(\to\)@
  (s' = s'') := by
  -- Induction over the shape
  induction h
  generalizing s'' with
        \end{Lean_Code}
\end{leftcolumn}
\begin{rightcolumn}
We assume that \(\langle S,\ s \rangle \to s'\) and shall prove that
if \(\langle S,\ s \rangle \to s''\) then \(s' = s''\).
We shall proceed by induction on the shape of the derivation tree for \(\langle S,\ s \rangle \to s'\).
\end{rightcolumn}
\vspace{2em}


\switchcolumn[1]*[]
\begin{leftcolumn}
    \begin{Lean_Code}
  | ass =>
    intro h_deriv_alt
    cases h_deriv_alt with
    | ass => rfl
    \end{Lean_Code}
\end{leftcolumn}
\begin{rightcolumn}
\textbf{The case} [ass$_{\text{ns}}$]: Then \(S\) is \(x := a\) and \(s'\) is \(s[x\mapsto\mathcal{A}[\![a]\!]s]\).
The only axiom or rule that could be used to give \( \langle x := a, s \rangle \to s'' \) is [ass$_{\text{ns}}$],
so it follows that \(s''\) must be \(s[x \mapsto \mathcal{A}[\![a]\!] s]\) and thereby \(s' = s''\).
\end{rightcolumn}
\vspace{2em}


\switchcolumn[1]*[]
\begin{leftcolumn}
    \begin{Lean_Code}
  | skip =>
    intro h_deriv_alt
    cases h_deriv_alt with
    | skip => rfl
    \end{Lean_Code}
\end{leftcolumn}
\begin{rightcolumn}
\textbf{The case} [skip$_{\text{ns}}$]: Analogous.
\end{rightcolumn}
\vspace{2em}


\switchcolumn[1]*[]
\begin{leftcolumn}
    \begin{Lean_Code}
  | comp
      h_left h_right
      ih_left ih_right =>
    intro h_deriv_alt
    cases h_deriv_alt with
    | comp
        h_alt_left
        h_alt_right =>
      have h_mid_eq :=
        ih_left _ h_alt_left
      subst h_mid_eq
      exact ih_right _
             h_alt_right
    \end{Lean_Code}
\end{leftcolumn}
\begin{rightcolumn}
\textbf{The case} [comp$_{\text{ns}}$]: Assume that
\[
\langle S_1;\ S_2,\ s \rangle \to s'
\]
holds because
\[
\langle S_1,\ s \rangle \to s_0\ \mbox{and}\ \langle S_2,\ s_0 \rangle \to s'
\]
for some \(s_0\). The only rule that could be applied to give \(\langle S_1;\ S_2,\ s \rangle \to s''\) is
[comp$_{\text{ns}}$], so there is a state \(s_1\) such that
\[
\langle S_1,\ s \rangle \to s_1\ \mbox{and}\ \langle S_2,\ s_1 \rangle \to s'
\]
The induction hypothesis can be applied to the premise \(\langle S_1,\ s \rangle \to s_0\) and from
\(\langle S_1,\ s \rangle \to s_1\) we get \(s_0 = s_1\). Similarly, the induction hypothesis can be applied
to the premise \(\langle S_2,\ s_0 \rangle \to s'\) and from \(\langle S_2,\ s \rangle \to s''\) we get \(s' = s''\) as required.
\end{rightcolumn}
\vspace{2em}


\switchcolumn[1]*[]
\begin{leftcolumn}
    \begin{Lean_Code}
  | if_true
      h_cond_true h_deriv_then
      ih_then =>
    intro h_deriv_alt
    cases h_deriv_alt with
    | if_true
        _ h_deriv_then_alt =>
      exact ih_then _
             h_deriv_then_alt

    | if_false
        h_cond_false_alt _ =>
      rw [h_cond_true] at
          h_cond_false_alt
      contradiction
    \end{Lean_Code}
\end{leftcolumn}
\begin{rightcolumn}
\textbf{The case} [if$_{\text{ns}}^{\ \text{tt}}$]: Assume that
\[
\langle \mathtt{if}\ b\ \mathtt{then}\ S_1\ \mathtt{else}\ S_2,\ s \rangle \to s'
\]
holds because
\[
\mathcal{B}[\![b]\!] = \mathbf{tt}\ \mbox{and}\ \langle S_1,\ s \rangle \to s'
\]
From \(\mathcal{B}[\![b]\!] = \mathbf{tt}\) we get that the only rule that could be applied to give the
alternative \(\langle \mathtt{if}\ b\ \mathtt{then}\ S_1\ \mathtt{else}\ S_2,\ s \rangle \to s''\) is 
[if$_{\text{ns}}^{\ \text{tt}}$]. So it must be the case that
\[
\langle S_1,\ s \rangle \to s'
\]
But then the induction hypothesis can be applied to the premise \(\langle S_1,\ s \rangle \to s'\)
and from \(\langle S_1,\ s \rangle \to s''\) we get \(s' = s''\).
\end{rightcolumn}
\vspace{2em}


\switchcolumn[1]*[]
\begin{leftcolumn}
    \begin{Lean_Code}
  | if_false
      h_cond_false h_deriv_else
      ih_else =>
    intro h_deriv_alt
    cases h_deriv_alt with
    | if_false
        _ h_deriv_else_alt =>
      exact ih_else _
             h_deriv_else_alt

    | if_true
        h_cond_true_alt _ =>
      rw [h_cond_false] at
          h_cond_true_alt
      contradiction
    \end{Lean_Code}
\end{leftcolumn}
\begin{rightcolumn}
\textbf{The case} [if$_{\text{ns}}^{\ \text{ff}}$]: Analogous.
\end{rightcolumn}
\vspace{2em}


\switchcolumn[1]*[]
\begin{leftcolumn}
    \begin{Lean_Code}
  | while_true
      h_cond_true h_deriv_body
      h_deriv_rest
      ih_while_body
      ih_while_rest =>
    intro h_deriv_alt
    cases h_deriv_alt with
    | while_false
        h_cond_false_alt =>
      rw [h_cond_true] at
          h_cond_false_alt
      contradiction

    | while_true
        _ h_deriv_body_alt
        h_deriv_rest_alt =>
      have h_state_eq :=
        ih_while_body _
        h_deriv_body_alt
      subst h_state_eq
      exact ih_while_rest _
             h_deriv_rest_alt
    \end{Lean_Code}
\end{leftcolumn}
\begin{rightcolumn}
\textbf{The case} [while$_{\text{ns}}^{\ \text{tt}}$]: Assume that
\[
\langle \mathtt{while}\ b\ \mathtt{do}\ S,\ s \rangle \to s'
\]
because
\[
\mathcal{B}[\![b]\!] = \mathbf{tt}, \ \langle S,\ s \rangle \to s_0\ \mbox{and}\ \langle \mathtt{while}\ b\ \mathtt{do}\ S,\ s_0 \rangle \to s'
\]
The only rule that could be applied to give \(\langle \mathtt{while}\ b\ \mathtt{do}\ S,\ s \rangle \to s''\)
is [while$_{\text{ns}}^{\ \text{tt}}$] because \(\mathcal{B}[\![b]\!] = \mathbf{tt}\), and this means that
\[
\langle S,\ s \rangle \to s_1\ \mbox{and}\ \langle \mathtt{while}\ b\ \mathtt{do}\ S,\ s_1 \rangle \to s''
\]
must hold for some \(s_1\). Again the induction hypothesis can be applied to the premise
\(\langle S,\ s \rangle \to s_0\), and from \(\langle S,\ s \rangle \to s_1\) we get \(s_0 = s_1\).
Thus we have
\[
\langle \mathtt{while}\ b\ \mathtt{do}\ S,\ s_0 \rangle \to s' \mbox{and}\ \langle \mathtt{while}\ b\ \mathtt{do}\ S,\ s_0 \rangle \to s''
\]
Since \(\langle \mathtt{while}\ b\ \mathtt{do}\ S,\ s_0 \rangle \to s'\) is a premise of [while$_{\text{ns}}^{\ \text{tt}}$],
we can apply the induction hypothesis to it. From \(\langle \mathtt{while}\ b\ \mathtt{do}\ S,\ s_0 \rangle \to s''\)
we therefore get \(s' = s''\) as required.
\end{rightcolumn}
\vspace{2em}


\switchcolumn[1]*[]
\begin{leftcolumn}
    \begin{Lean_Code}
  | while_false
      h_cond_false =>
    intro h_deriv_alt
    cases h_deriv_alt with
    | while_true
        h_cond_true_alt
        _ _ =>
      rw [h_cond_false] at
          h_cond_true_alt
      contradiction

    | while_false
        _ =>
      trivial
    \end{Lean_Code}
\end{leftcolumn}
\begin{rightcolumn}
\textbf{The case} [while$_{\text{ns}}^{\ \text{ff}}$]: Straightforward.
\end{rightcolumn}
\end{paracol}
\vspace{1em}

\subsubsection*{Comparison and Result}
The formalization of determinism highlights several key differences in how
"obvious" steps are handled:
\begin{itemize}
    \item \textbf{Inversion and Case Analysis}:
    The textbook frequently uses the phrase "the only rule that could be applied
    is...". In Lean, this is explicitly handled by the \texttt{cases} tactic
    on the second derivation (\texttt{h\_deriv\_alt}). This forces the user to
    explicitly handle (or dismiss via contradiction) cases that the textbook
    ignores.
    
    \item \textbf{Handling Side Conditions}:
    In the \texttt{if} and \texttt{while} cases, the textbook assumes that because
    $\mathcal{B}[\![b]\!]$ is true, only the "true" rule applies. In the mechanized
    proof, we must explicitly show that the "false" rule choice leads to a
    \texttt{contradiction} (e.g., \texttt{true = false}).
    
    \item \textbf{Explicit Substitutions}:
    When the induction hypothesis reveals that an intermediate state is equal
    ($s_0 = s_1$), the mechanized proof requires an explicit \texttt{subst}
    to update all occurrences of the variable in the goal. The textbook performs
    this substitution implicitly in its narrative flow.
    
    \item \textbf{State Generalization}:
    To prove the theorem, the variable $s''$ must be \\
    \texttt{generalized} in the induction.
    This ensures the IH is strong enough to talk about
    \textit{any} potential second derivation, not just one specific fixed $s''$.
    
    \item \textbf{Explication of Mutually Recursive Hypotheses}:
    In the \texttt{while\_true} case, the textbook treats the second application
    of the induction hypothesis to the "rest" of the loop as a standard step.
    In Lean, this corresponds to an IH generated by the inductive nature of
    the Big-Step semantics, making explicit that we are inducting on the
    \textit{height} or \textit{structure} of the derivation tree rather than
    just the syntax of the statement.
\end{itemize}


%
%
%
\section{Formalizing Structural Operational Semantics (Small-Step)}
\label{formal_ns_sos_equiv:sec:sos:formalization}
Similar to the way we define natural semantics we can define structural operational semantics,\missingcitation\todo{Cite the origin of Small-Step semantics if not already cited in previous sections.}
with the biggest difference being that not all statements will terminate to a final state after the first evaluation,but
instead might require multiple steps to reach a final step for the operation. This is most notable for composition rules.

\vspace{1em}
\begin{paracol}{2}
    \begin{leftcolumn}
        \centering \textbf{Mechanized (Lean 4)}
    \end{leftcolumn}
    \begin{rightcolumn}
        \centering \textbf{Pen-and-Paper (Textbook)} \\\noindent
    \end{rightcolumn}

    \switchcolumn[1]*[]
    \begin{leftcolumn}
        \begin{Lean_Code}
inductive Config where
  | step :
      Stmt  @\(\to\)@ State  @\(\to\)@ Config
  | final :
      State  @\(\to\)@ Config

set_option quotPrecheck
    false in
set_option hygiene false in
notation:40
    "@\color{purple}\(\langle\)@" S "," s "@\color{purple}\(\rangle\)@" " @\color{purple}\(\to_{sos}\)@ "
    "@\color{purple}\(\langle\)@" S' "," s':40 "@\color{purple}\(\rangle\)@" =>
    small_step S s
      (Config.step S' s')

set_option quotPrecheck
    false in
set_option hygiene false in
notation:40
    "@\color{purple}\(\langle\)@" S "," s "@\color{purple}\(\rangle\)@" " @\color{purple}\(\to_{sos}\)@ "
    s':40 =>
    small_step S s
      (Config.final s')
        \end{Lean_Code}
    \end{leftcolumn}
    \begin{rightcolumn}
    \end{rightcolumn}
    \vspace{2em}


    \switchcolumn[1]*[]
    \begin{leftcolumn}
        \begin{Lean_Code}
inductive small_step :
  Stmt @\(\to\)@ State @\(\to\)@
  Config @\(\to\)@ Prop where
        \end{Lean_Code}
    \end{leftcolumn}
    \begin{rightcolumn}
    \end{rightcolumn}
    \vspace{2em}


    \switchcolumn[1]*[]
    \begin{leftcolumn}
        \begin{Lean_Code}
  | ass {s x a} :
    @\(\langle\)@x @\(: \equiv\)@ a, s@\(\rangle\)@ @\(\to_{sos}\)@
      (s[x @\mapsto@ @\(\mathcal{A}[\![a]\!]\)@ s])
\end{Lean_Code}
    \end{leftcolumn}
    \begin{rightcolumn}
        \vspace{1.1em}
            \(
\begin{prooftree}
  \hypo{} % Empty hypothesis for an axiom
  \infer[
    left label=[ass$_{\text{sos}}$]
    ]1{
        \parbox{4cm}{$
    \langle x := a, s \rangle \Rightarrow \\
    s[x \mapsto \mathcal{A}[\![a]\!] s]
        $}
    }
\end{prooftree}
            \)
    \end{rightcolumn}
    \vspace{1em}


    \switchcolumn[1]*[]
    \begin{leftcolumn}
        \begin{Lean_Code}
  | skip {s} :
    @\(\langle\)@Stmt.skip, s@\(\rangle\)@ @\(\to_{sos}\)@ s
        \end{Lean_Code}
    \end{leftcolumn}
    \begin{rightcolumn}
        \vspace{1.1em}
            \(
\begin{prooftree}
  \hypo{} % Empty hypothesis for an axiom
  \infer[left label=[skip$_{\text{sos}}$]]1{ \langle \mathtt{skip}, s \rangle \Rightarrow s}
\end{prooftree}
            \)
    \end{rightcolumn}
    \vspace{1em}
    \switchcolumn[1]*[]
    \begin{leftcolumn}
        \begin{Lean_Code}
  | comp1 {S@$_{1}$@ S@$_{1}$@' S@$_{2}$@ s s'}
    (h_step :
    @\(\langle\)@S@$_{1}$@, s@\(\rangle\)@ @\(\to_{sos}\)@ @\(\langle\)@S@$_{1}$@', s'@\(\rangle\)@) :
    @\(\langle\)@S@$_{1}$@; S@$_{2}$@, s@\(\rangle\)@ @\(\to_{sos}\)@
      @\(\langle\)@S@$_{1}$@'; S@$_{2}$@, s'@\(\rangle\)@
        \end{Lean_Code}
    \end{leftcolumn}
    \begin{rightcolumn}
        \vspace{2.5em}
            \(
\begin{prooftree}
  \hypo{\langle S_1, s \rangle \Rightarrow \langle S_1', s' \rangle}
  \infer[
    left label=[comp$_{\text{sos}}^{1}$]
    ]1{
        \parbox{4cm}{$
    \langle S_1;\ S_2, s \rangle \Rightarrow \\
    \langle S_1';\ S_2, s' \rangle
        $}
    }
\end{prooftree}
            \)
    \end{rightcolumn}
    % \vspace{3em}

    \switchcolumn[1]*[]
    \begin{leftcolumn}
        \begin{Lean_Code}
  | comp2 {S@$_{1}$@ S@$_{2}$@ s s'}
    (h_left :
      @\(\langle\)@S@$_{1}$@, s@\(\rangle\)@ @\(\to_{sos}\)@ s') :
    @\(\langle\)@S@$_{1}$@; S@$_{2}$@, s@\(\rangle\)@ @\(\to_{sos}\)@ @\(\langle\)@S@$_{2}$@, s'@\(\rangle\)@
        \end{Lean_Code}
    \end{leftcolumn}
    \begin{rightcolumn}
        \vspace{2.5em}
            \(
\begin{prooftree}
  \hypo{\langle S_1, s \rangle \Rightarrow s'}
  \infer[
    left label=[comp$_{\text{sos}}^{2}$]
  ]1{
    \langle S_1;\ S_2, s \rangle \Rightarrow \langle S_2, s'\rangle}
\end{prooftree}
            \)
    \end{rightcolumn}
    % \vspace{3em}

    \switchcolumn[1]*[]
    \begin{leftcolumn}
        \begin{Lean_Code}
  | if_true {b S@$_{1}$@ S@$_{2}$@ s s'}
    (h_cond_true :
      @\(\mathcal{B}[\![b]\!]\)@ s = true) : 
    @\(\langle\)@if b then S@$_{1}$@ else S@$_{2}$@, s@\(\rangle\)@
      @\(\to_{sos}\)@ @\(\langle\)@S@$_{1}$@, s@\(\rangle\)@
        \end{Lean_Code}
    \end{leftcolumn}
    \begin{rightcolumn}
        \vspace{3.6em}
            \(
\begin{prooftree}
  \hypo{} % Empty for axiom
  \infer[
    left label={[if$_{\text{sos}}^{\text{\ tt}}$]}
    ]1[
    if $\mathcal{B}[\![b]\!]s = \mathbf{tt}$
    ]{
        \parbox{4cm}{$
    \langle \mathtt{if}\ b\ \mathtt{then}\ S_1\ \mathtt{else}\ S_2, s \rangle \\
    \Rightarrow \langle S_1, s\rangle
        $}
    }
\end{prooftree}
    \)
\end{rightcolumn}
% \vspace{3em}

\switchcolumn[1]*[]
\begin{leftcolumn}
    \begin{Lean_Code}
  | if_false {b S@$_{1}$@ S@$_{2}$@ s s'}
    (h_cond_false :
      @\(\mathcal{B}[\![b]\!]\)@ s = false) : 
    @\(\langle\)@if b then S@$_{1}$@ else S@$_{2}$@, s@\(\rangle\)@
      @\(\to_{sos}\)@ @\(\langle\)@S@$_{2}$@, s@\(\rangle\)@
    \end{Lean_Code}
\end{leftcolumn}
\begin{rightcolumn}
\vspace{3.6em}
    \(
\begin{prooftree}
  \hypo{} % Empty for axiom
  \infer[
    left label={[if$_{\text{sos}}^{\text{\ ff}}$]}
    ]1[
    if $\mathcal{B}[\![b]\!]s = \mathbf{ff}$
    ]{
        \parbox{4cm}{$
    \langle \mathtt{if}\ b\ \mathtt{then}\ S_1\ \mathtt{else}\ S_2, s \rangle \\
    \Rightarrow \langle S_2, s\rangle
        $}
    }
\end{prooftree}
            \)
    \end{rightcolumn}
    % \vspace{3em}

    \switchcolumn[1]*[]
    \begin{leftcolumn}
        \begin{Lean_Code}
  | while_unroll {b S s} :
    @\(\langle\)@while b then S, s@\(\rangle\)@ @\(\to_{sos}\)@
      @\(\langle\)@if b
      then (S; while b then S)
      else (Stmt.skip), s@\(\rangle\)@
    \end{Lean_Code}
\end{leftcolumn}
\begin{rightcolumn}
\vspace{1.1em}
    \(
\begin{prooftree}
  \hypo{}
  \infer[
      left label={[while$_{\text{sos}}$]}
      ]1{
        \parbox{4cm}{$
            \langle \mathtt{while}\ b\ \mathtt{do}\ S, s \rangle \Rightarrow \\
            \langle \mathtt{if}\ b\ \\
            \mathtt{then}\ (S;\ \mathtt{while}\ b\ \mathtt{do}\ S)\ \\
            \mathtt{else}\ \mathtt{skip}, s \rangle
        $}
    }
\end{prooftree}
    \)
\end{rightcolumn}
\end{paracol}
\vspace{1em}

\subsubsection*{Formalization of Structural Operational Semantics}

The transition system for Structural Operational Semantics is formalized
as an inductive relation \texttt{small\_step} that captures single-step
program execution.
The rules themselves closely mirror the textbook\missingcitation\todo{Cite textbook for the origin of these specific rules.}:
\begin{itemize}
    \item \textbf{Configuration Explicitly Typed}:
    To satisfy the Lean kernel, the implicit distinction between a "state"
    (termination) and a "statement-state pair" (continuation) must be made
    explicit via a \texttt{Config} inductive type. This addresses RQ2 by
    formalizing the "pedagogical shortcut" where textbook notation uses the
    same arrow for two different codomain types.
    
    \item \textbf{Type-System Driven Branching}:
    While the textbook presents [comp$^1$] and [comp$^2$] as visually similar
    rules, the mechanized version (RQ1) requires these to be mapped to the
    \texttt{step} and \texttt{final} constructors respectively. This structural
    reformulation ensures the Lean kernel can type-check the transition
    relation's codomain.

    \item \textbf{Conditional Rules}:
    The two conditional branches are explicitly checked by moving the side
    condition to be a premise instead, similarly to what was done for
    Natural Semantics.
\end{itemize}

While the single-step relation \texttt{small\_step} captures individual
program transitions, practical reasoning about complete program execution
requires multi-step relations. The textbook introduces convenient shorthand
notation for these ($\langle S, s\rangle \Rightarrow^{k} s'$ for exactly
$k$ steps and $\langle S, s\rangle \Rightarrow^{*}$ s' for any finite number
of steps)~\citep[p.~34 --- 35]{nielson_semantics_2007}. The mechanized version must explicitly formalize these as inductive
types to enable rigorous proof construction. To use the same ideas in Lean 4,
we introduce two new inductive types as follows:

\begin{Lean_Code}
inductive small_step_k : Config @\(\to\)@ Nat @\(\to\)@ Config @\(\to\)@ Prop where
  | refl {@\gamma@} :
      small_step_k @\gamma@ 0 @\gamma@
  | step {S s @\gamma@' @\gamma@'' k}
    (step : small_step S s @\gamma@')
    (rest : small_step_k @\gamma@' k @\gamma@'') :
      small_step_k (Config.step S s) (k + 1) @\gamma@''

-- Define the standard Nielson & Nielson notation
-- Non terminal configuration
notation:40 "@\color{purple}$\langle$@" S "," s "@\color{purple}$\rangle$@" @\leanbreak@ "@\ \color{purple}\to$_{sos}[$\ @" k "@\ \color{purple}$]$\ @" "@\color{purple}$\langle$@" S' "," s':40 "@\color{purple}$\rangle$@" => @\leanbreak@ small_step_k (Config.step S s) k (Config.step S' s')

-- Terminal configuration
notation:40 "@\color{purple}$\langle$@" S "," s "@\color{purple}$\rangle$@" @\leanbreak@ "@\ \color{purple}\to$_{sos}[$\ @" k "@\ \color{purple}$]$\ @" s':40 => @\leanbreak@ small_step_k (Config.step S s) k (Config.final s')

\end{Lean_Code}
and
\begin{Lean_Code}
inductive small_step_star : Config @\(\to\)@ Nat @\(\to\)@ Config @\(\to\)@ Prop where
  | refl {@\gamma@} :
      small_step_star @\gamma@ 0 @\gamma@
  | step {S s @\gamma@' @\gamma@'' k}
    (step : small_step S s @\gamma@')
    (rest : small_step_star @\gamma@' @\gamma@'') :
      small_step_star (Config.step S s) @\gamma@''

-- Define the standard Nielson & Nielson notation
-- Non terminal configuration
notation:40 "@\color{purple}$\langle$@" S "," s "@\color{purple}$\rangle$@" @\leanbreak@ "@\ \color{purple}\to$_{sos}*$\ @" "@\color{purple}$\langle$@" S' "," s':40 "@\color{purple}$\rangle$@" => @\leanbreak@ small_step_star (Config.step S s) (Config.step S' s')

-- Terminal configuration
notation:40 "@\color{purple}$\langle$@" S "," s "@\color{purple}$\rangle$@" @\leanbreak@ "@\ \color{purple}\to$_{sos}*$\ @" s':40 => @\leanbreak@ small_step_star (Config.step S s) (Config.final s')

\end{Lean_Code}

\paragraph{Note on Notation Scope:}
It should be noted that we deliberately do not introduce specific notation for
transitions starting from an abstract configuration
$\gamma$ (e.g., $\gamma \to_{sos} \gamma'$). While Lean's Config type encompasses
both step and final states, extending the custom notation to cover all permutations
of $\gamma$, Config.step, and Config.final would significantly complicate the
notation binding. Such a high degree of overloading often leads to ambiguity
in Lean's parser, causing more diagnostic issues than the shorthand would resolve.
Consequently, we formalize only the most common pedagogical mappings—from a
statement-state pair ($\langle S, s \rangle$) to either a successor configuration
or a terminal state—leaving transitions between abstract configurations to be
handled via manual applications of the underlying small\_step\_k and
small\_step\_star relations.

%
%
%
\section{Formalizing Proofs in Structural Operational Semantics}
Proofs made for structural operational semantics in the literature tend to share a pattern in using induction over the length of the derivation sequence.
More formally the book uses the following definition:
\begin{quotation}
\noindent
Induction on the Length of Derivation Sequences\vspace{1em} \\  \noindent
1: Prove that the property holds for all derivation sequences of length 0. \vspace{1em} \\ \noindent
2: Prove that the property holds for all finite derivation sequences:
Assume that the property holds for all derivation sequences of length at
most \(k\) (this is called the induction hypothesis) and show that it holds
for derivation sequences of length \(k+1\).
\end{quotation}
Doing induction over the length will be used in the proof in all structural operational semantics proofs in this thesis.
Let's start with a simpler proof before moving on to the final boss of equivalence between the two semantics.

%
%
\subsection{Decomposition (Lemma 2.19)}
\label{formal_ns_sos_equiv:sos:comp_split}
The first real proof provided for structural operational semantics is the proof that states
that a composition of two statements which terminates, can be split into two individual executions for each statement.

\begin{quotation}
    \noindent
    If \(\langle S_1;\ S_2, s \rangle \Rightarrow^k s''\), then there exists
    a state \(s'\) and natural numbers \(k_1\) and \(k_2\) such that
    \(\langle S_1, s \rangle \Rightarrow^{k_1} s'\) and \(\langle S_2, s' \rangle \Rightarrow^{k_2} s''\),
    where \(k = k_1 + k_2\).
\end{quotation}



\vspace{1em}
\begin{paracol}{2}
\begin{leftcolumn}
\centering \textbf{Mechanized (Lean 4)}
\end{leftcolumn}
\begin{rightcolumn}
\centering \textbf{Pen-and-Paper (Textbook)} \\\noindent
\end{rightcolumn}

\switchcolumn[1]*[]
\begin{leftcolumn}
    \begin{Lean_Code}
lemma seq_split {S@$_1$@ S@$_2$@ : Stmt}
    {s s'' : State} {k : Nat}
    (h : @\(\langle S_1; S_2, s \rangle \to_{sos}[k] s''\)@) :
    @\(\exists\)@ s' k@\(_1\)@ k@\(_2\)@, @\(\langle S_1, s \rangle \to_{sos}[k_1] s'\)@
    @\(\land\)@ @\(\langle S_2, s' \rangle \to_{sos}[k_2] s''\)@
    @\(\land\)@ k = k@\(_1\)@ + k@\(_2\)@ := by
  -- Proof by induction on
  -- the number k
  -- (length of the
  -- derivation sequence).
  induction k
  generalizing s S@$_1$@ with
    \end{Lean_Code}
\end{leftcolumn}
\begin{rightcolumn}
Proof: The proof is by induction on the number \(k\); that is, by induction on the
length of the derivation sequence \\
\(\langle S_1;\ S_2, s \rangle \Rightarrow^k s''\).
\end{rightcolumn}
\vspace{2em}

\switchcolumn[1]*[]
\begin{leftcolumn}
    \begin{Lean_Code}
  | zero =>
    -- Base case k = 0:
    -- impossible for a
    -- starting configuration
    -- @\color{green!40!black}{\(\langle S_1; S_2, s \rangle\)}@
    cases h
    \end{Lean_Code}
\end{leftcolumn}
\begin{rightcolumn}
If \(k = 0\), then the result holds vacuously
(because \(\langle S_1;\ S_2, s \rangle\) and \(s''\) are different).
\end{rightcolumn}
\vspace{2em}

\switchcolumn[1]*[]
\begin{leftcolumn}
    \begin{Lean_Code}
  | succ k@\color{green!40!black}{\(_0\)}@ ih =>
    -- Inductive step:
    -- assume the statement
    -- for @\color{green!40!black}{\(k_0\)}@,
    -- prove for k@\color{green!40!black}{\(_0\)}@ + 1.
    -- The derivation can be
    -- written as a first
    -- small-step to some
    -- configuration @\color{green!40!black}{\(\gamma\)}@,
    -- followed by k@\color{green!40!black}{\(_0\)}@ further
    -- steps to s''.
    cases h with
    | step gamma_step h_tail =>
    \end{Lean_Code}
\end{leftcolumn}
\begin{rightcolumn}
For the induction step, we assume that the lemma hold for \(k \leq k_0\), and we
shall prove it for \(k_0 + 1\). So assume that
\[
\langle S_1;\ S_2, s \rangle \Rightarrow^{k_0 + 1} s''
\]
This mean that the derivation sequence can be written as
\[
\langle S_1;\ S_2, s \rangle \Rightarrow \gamma \Rightarrow^{k_0} s''
\]
for some configuration \(\gamma\).
\end{rightcolumn}
\vspace{2em}

\switchcolumn[1]*[]
\begin{leftcolumn}
    \begin{Lean_Code}
      -- Now inspect which rule
      -- produced that first step:
      -- [comp1] or [comp2].
      cases gamma_step with
    \end{Lean_Code}
\end{leftcolumn}
\begin{rightcolumn}
Now one of two cases applies depending on
which of the two rules [comp\(_{sos}^{1}\)] and [comp\(_{sos}^{2}\)] is used.
\end{rightcolumn}
\vspace{2em}

\switchcolumn[1]*[]
\begin{leftcolumn}
    \begin{Lean_Code}
      | comp1 progress =>
        -- [comp1]:
        -- @\color{green!40!black}{\(\gamma = \langle S_1'; S_2, s' \rangle\)}@
        -- because progress :
        -- @\color{green!40!black}{\(\langle S_1, s \rangle \to_{sos} \langle S_1', s' \rangle\)}@.
        -- We therefore have
        -- @\color{green!40!black}{\(\langle S_1'; S_2, s' \rangle \to_{sos}[k_0] s''\)}@,
        -- so apply IH to that
        -- shorter derivation.
        specialize ih h_tail
        let
          @\(\langle s_0, k_1, k_2, hk_1, hk_2, h_{sum} \rangle\)@
          := ih
    \end{Lean_Code}
\end{leftcolumn}
\begin{rightcolumn}
In the first case, were [comp\(_{sos}^{1}\)] is used, we have
\(\gamma = \langle S_1';\ S_2, s' \rangle\) and
\[
\langle S_1;\ S_2, s \rangle \Rightarrow \langle S_1';\ S_2, s' \rangle
\]
because
\[
\langle S_1, s \rangle \Rightarrow \langle S_1', s' \rangle
\]
We therefore have
\[
\langle S_1';\ S_2, s' \rangle \Rightarrow^{k_0} s''
\]
and the induction hypothesis can be applied to this derivation sequence because it
is shorter than the one with which we started.
This means that there is a state \(s_0\) and natural numbers \(k_1\) and \(k_2\) such that
\[
\langle S_1', s' \rangle \Rightarrow^{k_1} s_0\ \mbox{and}\ \langle S_2, s_0 \rangle \Rightarrow^{k_2} s''
\]
where \(k_1 + k_2 = k_0\).
\end{rightcolumn}
\vspace{2em}

\switchcolumn[1]*[]
\begin{leftcolumn}
    \begin{Lean_Code}
        -- From progress:
        -- @\color{green!40!black}{\(\langle S_1, s \rangle \to_{sos} \langle S_1', s' \rangle\)}@
        -- and hk@\color{green!40!black}{\(_1\)}@ :
        -- @\color{green!40!black}{\(\langle S_1', s' \rangle \to_{sos}[k_1]\ s_0\)}@
        -- we obtain
        -- @\color{green!40!black}{\(\langle S_1, s \rangle \to_{sos}[k_1 + 1]\ s_0\)}@
        exists s@\(_0\)@, k@\(_1\)@ + 1, k@\(_2\)@

        constructor
        case left =>
          exact small_step_k.step
                progress
                hk@\(_1\)@

        case right =>
          constructor
          case left => exact hk@\(_2\)@

          case right =>
            -- arithmetic:
            -- (k@\color{green!40!black}{\(_1\)}@ + 1) + k@\color{green!40!black}{\(_2\)}@ =
            -- k@\color{green!40!black}{\(_0\)}@ + 1
            simp [h_sum]
            linarity
    \end{Lean_Code}
\end{leftcolumn}
\begin{rightcolumn}
Using that \(\langle S_1, s \rangle \Rightarrow \langle S_1', s' \rangle\) and \\
\(\langle S_1', s' \rangle \Rightarrow^{k_1} s_0\), we get
\[
\langle S_1, s \rangle \Rightarrow^{k_1+1} s_0
\]
We have already seen that \(\langle S_2, s_0 \rangle \Rightarrow^{k_2} s''\), and since
\((k_1 + 1) + k_2 = k_0 + 1\) we have proved the required result.
\end{rightcolumn}
\vspace{2em}

\switchcolumn[1]*[]
\begin{leftcolumn}
    \begin{Lean_Code}
      | comp2 terminates =>
        -- [comp2]:
        -- the first step
        -- returns @\color{green!40!black}{\(\gamma = \langle S_2, s' \rangle\)}@
        -- because
        -- terminates :
        -- @\color{green!40!black}{\(\langle S_1, s \rangle \to_{sos} s'\)}@.
        -- We have
        -- @\color{green!40!black}{\(\langle S_2, s' \rangle \to_{sos}[k_0]\ s''\)}@.
        -- Choose k@\color{green!40!black}{\(_1\)}@ = 1
        -- and k@\color{green!40!black}{\(_2\)}@ = k@\color{green!40!black}{\(_0\)}@.
        rename_i s'
        exists s', 1, k@\(_0\)@
        constructor
        case left =>
          exact small_step_k.step terminates small_step_k.refl

        case right =>
          constructor
          case left =>
            exact h_tail
          
          case right => linarith
    \end{Lean_Code}
\end{leftcolumn}
\begin{rightcolumn}
The second possibility is that [comp\(_{sos}^{1}\)]
has been used to obtain the derivation \\
\(\langle S_1;\ S_2, s \rangle \Rightarrow \gamma\)
Then we have
\[
\langle S_1, s \rangle \Rightarrow s'
\]
and \(\gamma\) is \(\langle S_2, s' \rangle\) so that
\[
\langle S_2, s' \rangle \Rightarrow^{k_0} s''
\]
The result now follows by choosing \(k_1 = 1\) and \(k_2 = k_0\).
\end{rightcolumn}


\end{paracol}
\vspace{1em}

\subsubsection*{Comparison and Result}
While the pen-and-paper proof follows a high-level logical flow, the mechanized
proof requires explicit handling of Lean-specific concepts and the formalization
of "pedagogical shortcuts":

\begin{itemize}
    \item \textbf{Formalization of the Relation}:
    The textbook uses a somewhat informal "where $k = k_1 + k_2$" clause.
    In Lean, this must be part of the explicit existential bundle ($\land k = k_1 + k_2$),
    requiring the use of \texttt{constructor} to split the logical conjunctions and
    \texttt{exists} to provide witnesses.
    
    \item \textbf{Variable Generalization}:
    The \texttt{generalizing s S$_1$} keyword is a crucial pedagogical gap.
    The textbook assumes the property holds for "all derivation sequences,"
    but in a tactic-based proof, one must explicitly ensure that the induction
    hypothesis is not bound to the specific initial state \(s\), or the proof
    will fail when the state changes during a step.

    \item \textbf{Explicit Case Exhaustion}:
    The textbook notes the \(k=0\) case is
    "vacuous" because the configurations are different. In Lean, this requires the
    \texttt{cases} tactic to demonstrate that no constructor in the inductive definition
    of the transition relation can produce a state from a statement in zero steps.

    \item \textbf{Specialization and Unpacking}:
    Tactics like \texttt{specialize}
    and the use of angle-bracket pattern matching (\(\langle s_0, k_1, \dots \rangle\))
    are required to bridge the gap between a general hypothesis and the specific
    sub-derivation (\texttt{h\_tail}) generated by the \texttt{step} constructor.

    \item \textbf{Arithmetic Bookkeeping}:
    The mechanized proof relies on \texttt{linarith}
    to solve the summation of steps (\((k_1 + 1) + k_2 = k_0 + 1\)). In the literature,
    this is dismissed as "the required result," but the Lean kernel requires a formal
    proof of equality between natural numbers.

    \item \textbf{Intermediate Configuration Inspection}:
    Lean requires an explicit
    second level of case analysis (\texttt{cases gamma\_step}) to identify which
    specific SOS rule was used to progress from \(\gamma\). The textbook performs
    this move implicitly by immediately discussing [comp\(_{sos}^{1}\)] and [comp\(_{sos}^{2}\)].
\end{itemize}


%
%
%
\section{Proving Natural Semantics $\equiv$ Structural Operational Semantics}
\label{formal_ns_sos_equiv:sec:ns_equiv_sos}

\todo{Fix an intro message, present proof, and why it is interesting to evaluate in thesis}
\todo{Explain that semantic function can be model in lean as it is partial, and one can not create
proofs using partial functions in lean 4}

\begin{quotation}
\noindent
For every statement S of \While, we have
\(\mathcal{S}_{ns}[\![S]\!] = \mathcal{S}_{sos}[\![S]\!]\).
This result expresses two properties:
\begin{itemize}
    \item If the execution of \(S\) from some state terminates in of of the semantics,
    then it also terminates in the other and the resulting states will be equal.
    \item If the execution of \(S\) from some state loops in one of the semantics,
    then it will also loop in the other.
\end{itemize}
It should be fairly obvious that the first property follows from the theorem because
there are no stuck configurations in the structural operational semantics of \While. \missingcitation \todo{Citation needed for the claim that there are no stuck configurations in this specific SOS.}
For the other property, suppose that the execution of \(S\) on state \(s\) loops in one semantics.
If it terminates in the other semantics, we have a contradiction with the first property
because both semantics are deterministic (Theorem 2.9 and Exercise 2.22). \missingcitation \todo{Cite the specific book where Theorem 2.9 and Exercise 2.22 are located.}
Hence \(S\) will have to loop on state \(s\) also in the other semantics.

The theorem is proved in two stages, as expressed by Lemma 2.27 and Lemma 2.28. \missingcitation \todo{Cite the source for Lemma 2.27 and 2.28.}
We shall first prove Lemma 2.27.~\citep[p.~41 --- 42]{nielson_semantics_2007}
\end{quotation}

\subsubsection*{Lemma 2.27}
For every statement \(S\) of \While and states \(s\) and \(s'\) we have
\[
\langle S, s \rangle \to s'\ \mbox{implies}\ \langle S, s \rangle \Rightarrow^{*} s'
\]
So if the execution of \(S\) from \(s\) terminates in the natural semantics, then it will terminate
in the same state in the structural operational semantics.
\vspace{1em}
\begin{paracol}{2}
\begin{leftcolumn}
\centering \textbf{Mechanized (Lean 4)}
\end{leftcolumn}
\begin{rightcolumn}
\centering \textbf{Pen-and-Paper (Textbook)}
\end{rightcolumn}

\switchcolumn[1]*[]
\begin{leftcolumn}
    \begin{Lean_Code}
theorem ns_to_sos
  (S : Stmt)
  (s s' : State) :
  ((@\(\langle S, s\rangle \to_{ns} s'\)@) @\(\to\)@
   (@\(\langle S, s \rangle \to_{sos}^{*} s'\)@))
  := by
  intro h
  induction h with
    \end{Lean_Code}
\end{leftcolumn}
\begin{rightcolumn} \noindent
Proof: The proof proceeds by induction on the shape of the derivation tree for\\
\(\langle S, s \rangle \to s'\).

\end{rightcolumn}
\vspace{2em}

\switchcolumn[1]*[]
\begin{leftcolumn}
    \begin{Lean_Code}
  | ass =>
    apply small_step_star.step
    apply small_step.ass
    apply small_step_star.refl
    \end{Lean_Code}
\end{leftcolumn}
\begin{rightcolumn}
\textbf{The case} [ass$_{\text{ns}}$]: We assume that
\[
\langle x:=a, s \rangle \to s[x \mapsto \mathcal{A}[\![a]\!] s]
\]
From [ass$_{\text{sos}}$], we get the required
\[
\langle x:=a, s \rangle \Rightarrow s[x \mapsto \mathcal{A}[\![a]\!] s]
\]
\end{rightcolumn}
\vspace{2em}


\switchcolumn[1]*[]
\begin{leftcolumn}
    \begin{Lean_Code}
  | skip =>
    apply small_step_star.step
    apply small_step.skip
    apply small_step_star.refl
    \end{Lean_Code}
\end{leftcolumn}
\begin{rightcolumn}
\textbf{The case} [skip$_{\text{ns}}$]: Analogous.
\end{rightcolumn}
\vspace{2em}


\switchcolumn[1]*[]
\begin{leftcolumn}
    \begin{Lean_Code}
  | comp ih1 ih2 h1 h2 =>
    rename_i S@\(_1\)@ S@\(_2\)@ s s' s''
    -- Exercise 2.21
    have h_comp_first_part :=
      exercise_2_21
        (S@\(_2\)@ := S@\(_2\)@) h1
    -- Transitivity
    exact
      small_step_star_trans
      h_comp_first_part h2
    \end{Lean_Code}
\end{leftcolumn}
\begin{rightcolumn}
\textbf{The case} [comp$_{\text{ns}}$]: Assume that
\[
\langle S_1;\ S_2,\ s \rangle \to s''
\]
because
\[
\langle S_1,\ s \rangle \to s'\ \mbox{and}\ \langle S_2,\ s' \rangle \to s''
\]
The induction hypothesis can be applied to both of the premises
\(\langle S_1,\ s \rangle \to s'\) and \(\langle S_2,\ s' \rangle \to s''\)
and gives
\[
\langle S_1,\ s \rangle \Rightarrow^{*} s'\ \mbox{and}\ \langle S_2,\ s' \rangle \Rightarrow^{*} s''
\]
From Exercise 2.21, we get
\[
\langle S_1;\ S_2,\ s \rangle \Rightarrow^{*} \langle S_2,\ s' \rangle
\]
and thereby \(\langle S_1;\ S_2,\ s \rangle \Rightarrow^{*} s''\).

\end{rightcolumn}
\vspace{2em}


\switchcolumn[1]*[]
\begin{leftcolumn}
    \begin{Lean_Code}
  | if_true h_cond h ih =>
    apply small_step_star.step
    · apply
        small_step.if_true
          h_cond
    · exact ih
    \end{Lean_Code}
\end{leftcolumn}
\begin{rightcolumn}
\textbf{The case} [if$_{\text{ns}}^{\ \text{tt}}$]: Assume that
\[
\langle \mathtt{if}\ b\ \mathtt{then}\ S_1\ \mathtt{else}\ S_2,\ s \rangle \to s'
\]
because
\[
\mathcal{B}[\![b]\!] = \mathbf{tt}\ \mbox{and}\ \langle S_1,\ s \rangle \to s'
\]
Since \(\mathcal{B}[\![b]\!] = \mathbf{tt}\), we get
\[
\langle \mathtt{if}\ b\ \mathtt{then}\ S_1\ \mathtt{else}\ S_2,\ s \rangle \Rightarrow \langle S_1,\ s \rangle \Rightarrow^{*} s'
\]
where the first relationship comes from [if$_{\text{ns}}^{\ \text{tt}}$] and the second
from the induction hypothesis applied to the premise\\
\(\langle S_1,\ s \rangle \to s'\).
\end{rightcolumn}
\vspace{2em}


\switchcolumn[1]*[]
\begin{leftcolumn}
    \begin{Lean_Code}
  | if_false h_cond h ih =>
    -- [if_false_sos]
    apply small_step_star.step
    · apply small_step.if_false h_cond
    · exact ih
    \end{Lean_Code}
\end{leftcolumn}
\begin{rightcolumn}
\textbf{The case} [if$_{\text{ns}}^{\ \text{ff}}$]: Analogous.
\end{rightcolumn}
\vspace{2em}


\switchcolumn[1]*[]
\begin{leftcolumn}
    \begin{Lean_Code}
  | while_true h_cond h_body h_rest ih_body ih_loop =>
    rename_i b S s s' s''
    -- Exercise 2.21
    have h_after_ex21 := exercise_2_21 (S@\(_2\)@ := Stmt.loop b S) ih_body
    -- while_sos step
    apply small_step_star.step
    · exact small_step.while_unroll
    -- if_sos^tt step
    apply small_step_star.step
    · exact small_step.if_true h_cond
    -- Transitivity
    exact small_step_star_trans h_after_ex21 ih_loop
    \end{Lean_Code}
\end{leftcolumn}
\begin{rightcolumn}
\textbf{The case} [while$_{\text{ns}}^{\ \text{tt}}$]: Assume that
\[
\langle \mathtt{while}\ b\ \mathtt{do}\ S,\ s \rangle \to s''
\]
because
\[
\mathcal{B}[\![b]\!] = \mathbf{tt}, \ \langle S,\ s \rangle \to s'\ \mbox{and}\ \langle \mathtt{while}\ b\ \mathtt{do}\ S,\ s' \rangle \to s''
\]
The induction hypothesis can be applied to both of the premises \(\langle S,\ s \rangle \to s'\)
and \(\langle \mathtt{while}\ b\ \mathtt{do}\ S,\ s' \rangle \to s''\)
and gives
\[
\langle S,\ s \rangle \Rightarrow^{*} s'\ \mbox{and}\ \langle \mathtt{while}\ b\ \mathtt{do}\ S,\ s' \rangle \Rightarrow^{*} s''
\]
Using Exercise 2.21, we get
\[
\langle S;\ \mathtt{while}\ b\ \mathtt{do}\ S,\ s \rangle \Rightarrow^{*} s''
\]
Using [while$_{\text{sos}}$] and [if$_{\text{sos}}^{\ \text{tt}}$] (with \(\mathcal{B}[\![b]\!] = \mathbf{tt}\)),
we get the first two steps of
\begin{align*}
    &\langle \mathtt{while}\ b\ \mathtt{do}\ S,\ s \rangle \\ & \Rightarrow
    \langle \mathtt{if}\ b\ \mathtt{then}\ \left( S;\ \mathtt{while}\ b\ \mathtt{do}\ S \right)\ \mathtt{else}\ \mathtt{skip},\ s \rangle \\ & \Rightarrow
    \langle S;\ \mathtt{while}\ b\ \mathtt{do}\ S,\ s \rangle \\ &\Rightarrow^{*}
    s''
\end{align*}
and we have already argued for the last part.
\end{rightcolumn}
\vspace{2em}


\switchcolumn[1]*[]
\begin{leftcolumn}
    \begin{Lean_Code}
  | while_false h_cond =>
    -- while_sos step
    apply small_step_star.step
    · exact small_step.while_unroll
    -- if_sos^ff step
    apply small_step_star.step
    · exact small_step.if_false h_cond
    -- skip_sos step + reflexivity
    apply small_step_star.step
    · exact small_step.skip
    · exact small_step_star.refl
    \end{Lean_Code}
\end{leftcolumn}
\begin{rightcolumn}
\textbf{The case} [while$_{\text{ns}}^{\ \text{ff}}$]: Straightforward.
\end{rightcolumn}
\end{paracol}
\vspace{1em}
This completes the proof of Lemma 2.27. The second part of the theorem
follows from Lemma 2.28.

\subsubsection*{Comparison and Result (Lemma 2.27)}
The mechanized proof of ns\_to\_sos mirrors the textbook's inductive structure
but requires explicit handling of Lean-specific machinery and the resolution
of pedagogical shortcuts:

\begin{itemize}
    \item \textbf{Constructor-Based Induction}:
    The textbook performs induction on "the shape of the derivation tree" informally.
    Lean explicitly matches on the constructors of the inductive type representing
    the natural semantics (\texttt{ass}, \texttt{skip}, \texttt{comp}, etc.),
    ensuring every possible derivation tree is covered
    
    \item \textbf{Implicit vs. Explicit Multi-stepping}:
    A significant "pedagogical shortcut" in the textbook is the conflation of
    single steps ($\Rightarrow$) and multi-steps ($\Rightarrow^*$).
    Lean requires explicit use of \texttt{small\_step\_star.step} to bridge this gap,
    demanding a formal "count" of steps that the textbook elides with phrases
    like "we get the required..."
    
    \item \textbf{Transitivity and Composition}:
    The textbook's logical chaining (e.g., "we have already argued for the last part")
    is formalized via \texttt{small\_step\_star\_trans}. This requires the user to provide
    intermediate states ($s'$) that are often left implicit in written proofs
    
    \item \textbf{Formalization of "Analogous" and "Straightforward"}:
    The textbook uses these terms to skip the \texttt{skip}, \texttt{if\_false}, and
    \texttt{while\_false} cases. In Lean, these require full tactic blocks.
    Even "straightforward" cases like \texttt{while\_false} require a three-step sequence
    (\texttt{while\_unroll} $\to$ \texttt{if\_false} $\to$ \texttt{skip}) that the
    literature treats as a single atomic intuition
    
    \item \textbf{Lemma Injection}:
    The textbook references Exercise 2.21 as a contextual hint.
    In Lean, this is a structural necessity; the proof cannot progress without an
    explicit application of \texttt{exercise\_2\_21}, often requiring manual
    instantiation of variables (e.g., \texttt{S2 := Stmt.loop b S}) to help
    the unifier
    
    \item \textbf{Variable Scoping}:
    The use of \texttt{rename\_i} highlights a mapping difference: textbooks use
    global naming conventions for states ($s, s', s''$), whereas Lean generates
    internal names during induction that must be explicitly mapped to the thesis's
    notation for readability.
\end{itemize}


\subsubsection*{Lemma 2.28}
For every statement \(S\) of While, states \(s\) and \(s'\) and natural number \(k\), we
have that
\[
\langle S, s \rangle \Rightarrow^{k} s'\ \mbox{implies}\ \langle S, s \rangle \to s'
\]
So if the execution of \(S\) from \(s\) terminates in the structural operational semantics,
then it will terminate in the same state in the natural semantics.



\vspace{1em}
\begin{paracol}{2}
\begin{leftcolumn}
\centering \textbf{Mechanized (Lean 4)}
\end{leftcolumn}
\begin{rightcolumn}
\centering \textbf{Pen-and-Paper (Textbook)}
\end{rightcolumn}

\switchcolumn[1]*[]
\begin{leftcolumn}
    \begin{Lean_Code}
lemma sos_k_to_ns (S : Stmt) (s s' : State) (k : Nat) :
  (@\(\langle\)@S, s@\(\rangle\)@ @\(\to_{sos}\)@[k] s') @\(\to\)@ @\(\mathbb{Z}\)@@\(\langle\)@S, s@\(\rangle\)@ @\(\to_{ns}\)@ s') := by
  -- Use strong induction on k to allow for arbitrary k1, k2 < k
  induction k using Nat.strongRecOn generalizing S s s'
  rename_i n ih
  intro h_k_step
  match n with
    | 0 =>
      -- Case k = 0: Vacuous as per text
      cases h_k_step
    | k@\(_0\)@ + 1 =>
      -- Case k = k0 + 1: This is where we "inspect the first step"
      cases h_k_step with
      | step h_first h_rest =>
        -- h_first: @\(\langle\)@S, s@\(\rangle\)@ @\(\to_{sos}\)@ next_cfg
        -- h_rest : next_cfg @\(\to_{sos}\)@[k0] s'

        -- Now you can perform `cases h_first` to inspect the SOS rule tree
        cases h_first
    \end{Lean_Code}
\end{leftcolumn}
\begin{rightcolumn} \noindent
Proof: The proof proceeds by induction on the length of of the derivation sequence 
\(\langle S, s \rangle \Rightarrow^{k} s'\); that is, by induction on \(k\)

If \(k = 0\), then the result holds vacuously. \\


To prove the the induction step we assume that the lemma holds for \(k \leq k_0\),
and we shall then prove that it holds for \(k_0 + 1\).
We proceed by cases on how the first step of \(\langle S, s \rangle \Rightarrow^{k_+1} s'\)
is obtained; that is, by inspecting the derivation tree for the first step of computation in structural operational semantics.
\end{rightcolumn}
\vspace{2em}

\switchcolumn[1]*[]
\begin{leftcolumn}
    \begin{Lean_Code}
        case ass x a =>
          -- Text: "The case [ass_sos]: Straightforward (and k0 = 0)"
          cases h_rest -- This must be refl because assignment terminates in 1 step
          apply big_step.ass
    \end{Lean_Code}
\end{leftcolumn}
\begin{rightcolumn}
\textbf{The case} [ass$_{\text{sos}}$]: Straightforward (and \(k_0 = 0\)).
\end{rightcolumn}
\vspace{2em}


\switchcolumn[1]*[]
\begin{leftcolumn}
    \begin{Lean_Code}
        case skip =>
          -- Text: "The case [skip_sos]: Straightforward (and k0 = 0)"
          cases h_rest
          apply big_step.skip
    \end{Lean_Code}
\end{leftcolumn}
\begin{rightcolumn}
\textbf{The case} [skip$_{\text{sos}}$]: Straightforward (and \(k_0 = 0\)).
\end{rightcolumn}
\vspace{2em}


\switchcolumn[1]*[]
\begin{leftcolumn}
    \begin{Lean_Code}
        case comp1 S1 S1' S2 s_next h_step_S1 =>
          -- [Step 1: Apply Lemma 2.19 to split the sequence]
          obtain @\(\langle\)@s_mid, k@\(_1\)@, k@\(_2\)@, hk1, hk2, h_sum@\(\rangle\)@ := seq_split h_rest

          -- [Step 4: Use the [comp_ns] rule to combine results]
          apply big_step.comp (s' := s_mid)
          case h_left =>
            -- [Step 2: Establish the derivation for S1]
            have h_s1_full : @\(\langle\)@S1, s@\(\rangle\)@ @\(\to_{sos}\)@[k@\(_1\)@ + 1] s_mid := small_step_k.step h_step_S1 hk1
            -- [Step 3a: Justify that k@\(_1\)@ + 1 is strictly smaller than the total]
            have h_k2_pos : k@\(_2\)@ > 0 := by
              cases k@\(_2\)@
              case zero => cases hk2
              case succ => linarith
            -- [Step 3b: Apply the Induction Hypothesis (IH)]
            apply ih (k@\(_1\)@ + 1) (by linarith) S1 s s_mid h_s1_full
            
          case h_right =>
            -- [Step 3: Apply the Induction Hypothesis (IH) for S2]
            apply ih k@\(_2\)@ (by linarith) S2 s_mid s' hk2

        case comp2 S1 S2 s_next h_step_s1 =>
          apply big_step.comp (s' := s_next)
          case h_left =>
            -- Unlike comp1, S1 terminates fully in this first step (S1' is skip)
            -- We apply the IH to a 1-step derivation instead of (k@\(_1\)@ + 1)
            have h_lt_one : 1 < k@\(_0\)@ + 1 := by
              cases k@\(_0\)@
              case zero => cases h_rest
              case succ => linarith

            apply ih 1 h_lt_one S1 s s_next
            exact small_step_k.step h_step_s1 small_step_k.refl

          case h_right =>
            -- In comp1, the second part takes k@\(_2\)@ steps; here, it takes the full remaining k@\(_0\)@ steps
            apply ih k@\(_0\)@ (by linarith) S2 s_next s' h_rest
    \end{Lean_Code}
\end{leftcolumn}
\begin{rightcolumn}
\textbf{The case} [comp$_{\text{sos}}^{\text{1}}$] and [comp$_{\text{sos}}^{\text{2}}$]:
In both cases we assume that
\[
\langle S_1;\ S_2,\ s \rangle \Rightarrow^{k_0 + 1} s''
\]
We can now apply Lemma 2.19 and get that there exist a state \(s'\) and natural numbers
\(k_1\) and \(k_2\) such that
\[
\langle S_1,\ s \rangle \Rightarrow^{k_1} s'\ \mbox{and}\ \langle S_2,\ s' \rangle \Rightarrow^{k_2} s''
\]
where \(k_1 + k_2 = k_0 + 1\). The induction hypothesis can now be applied to each of these
derivation sequences because \(k_1 \leq k_0\) and \(k_2 \leq k_0\). So we get
\[
\langle S_1,\ s \rangle \to s'\ \mbox{and}\ \langle S_2,\ s' \rangle \to s''
\]
Using [comp$_{\text{ns}}$], we now get required \(\langle S_1;\ S_2,\ s \rangle \to s''\)
\end{rightcolumn}
\vspace{2em}


\switchcolumn[1]*[]
\begin{leftcolumn}
    \begin{Lean_Code}
        case if_true
            b S1 S2 h_cond =>
          apply big_step.if_true
            h_cond
          -- apply induction hypothesis
          apply ih k@\(_0\)@
            (by linarith) S1
            s s' h_rest
    \end{Lean_Code}
\end{leftcolumn}
\begin{rightcolumn}
\textbf{The case} [if$_{\text{sos}}^{\ \text{tt}}$]: Assume that
\(\mathcal{B}[\![b]\!] = \mathbf{tt}\) and that
\[
\langle \mathtt{if}\ b\ \mathtt{then}\ S_1\ \mathtt{else}\ S_2,\ s \rangle \Rightarrow \langle S_1,\ s \rangle \Rightarrow^{k_0} s'
\]
The induction hypothesis can be applied to the derivation \(\langle S_1,\ s \rangle \Rightarrow^{k_0} s'\) and
give
\[
\langle S_1,\ s \rangle \to s'
\]
The result now follows using [if$_{\text{ns}}^{\ \text{tt}}$].

\end{rightcolumn}
\vspace{2em}


\switchcolumn[1]*[]
\begin{leftcolumn}
    \begin{Lean_Code}
        case if_false
            b S1 S2 h_cond =>
          apply big_step.if_false
            h_cond
          apply ih k@\(_0\)@
            (by linarith) S2
            s s' h_rest
    \end{Lean_Code}
\end{leftcolumn}
\begin{rightcolumn}
\textbf{The case} [if$_{\text{sos}}^{\ \text{ff}}$]: Analogous.
\end{rightcolumn}
\vspace{2em}


\switchcolumn[1]*[]
\begin{leftcolumn}
    \begin{Lean_Code}
        case while_unroll b S_body =>
          -- First step: @\(\langle\)@while...@\(\rangle\)@ @\(\to\)@ @\(\mathbb{Z}\)@if b then (S; while...) else skip, s@\(\rangle\)@
          -- We apply the IH to the resulting k0 steps starting from the 'if'
          have h_ns_if := ih k@\(_0\)@ (by linarith) _ _ _ h_rest
          -- Then use the bridge: if_ns @\(\to\)@ @\(\mathbb{Z}\)@hile_ns (via lemma_2_5.mp)
          apply (lemma_2_5 b S_body s s').mp
          exact h_ns_if
    \end{Lean_Code}
\end{leftcolumn}
\begin{rightcolumn}
\textbf{The case} [while$_{\text{sos}}$]: We have

\begin{align*}
    &\langle \mathtt{while}\ b\ \mathtt{do}\ S,\ s \rangle \\ & \Rightarrow
    \langle \mathtt{if}\ b\ \mathtt{then}\ \left( S;\ \mathtt{while}\ b\ \mathtt{do}\ S \right)\ \mathtt{else}\ \mathtt{skip},\ s \rangle \\ & \Rightarrow
    \langle S;\ \mathtt{while}\ b\ \mathtt{do}\ S,\ s \rangle \\ &\Rightarrow^{k_0}
    s''
\end{align*}
The induction hypothesis can be applied to the \(k_0\) last steps of the derivation sequence and gives
\[
\langle \mathtt{if}\ b\ \mathtt{then}\ \left( S;\ \mathtt{while}\ b\ \mathtt{do}\ S \right)\ \mathtt{else}\ \mathtt{skip},\ s \rangle \to s''
\]
and from Lemma 2.5 we get the required
\[
\langle \mathtt{while}\ b\ \mathtt{do}\ S,\ s \rangle \to s''
\]
\end{rightcolumn}
\end{paracol}
\vspace{1em}

\subsubsection*{Comparison and Result (Lemma 2.28)}
The mechanized proof of sos\_k\_to\_ns demonstrates how to formalize a proof
by induction on the "length" of a derivation:

\begin{itemize}
    \item \textbf{Strong Induction (RQ1):}
    While the textbook specifies induction on $k$, Lean requires
    \texttt{Nat.strongRecOn}.This is a structural reformulation where the
    induction hypothesis must cover all $m < k$ to account for sequence
    splitting (Lemma 2.19) where $k_1$ and $k_2$ are arbitrary values
    summing to the total.
    
    \item \textbf{Explict Witness Handling (RQ2):}
    In the [comp] case, the textbook assumes the existence of an intermediate
    state $s'$ and step counts $k_1, k_2$. In Lean, these are obtained
    via the \texttt{obtain} tactic from a helper lemma, making the existential
    nature of the sequence split explicit to the kernel.
    
    \item \textbf{Termination Constraints (RQ2):}
    The textbook notes that for \texttt{skip} and \texttt{ass}, $k_0 = 0$ is
    "straightforward." Lean forces the user to prove this by showing that the
    "remaining steps" (\texttt{h\_rest}) in a derivation of a terminal command
    must be the empty/reflexive sequence.
    
    \item \textbf{Arithmetic Justification (RQ2):}
    The textbook glosses over the fact that if a sequence is split, the parts
    are strictly smaller than the whole. Lean requires explicit arithmetic
    proofs (using \texttt{linarith}) to satisfy the termination checker when
    applying the Induction Hypothesis.
    
    \item \textbf{Bridging Lemmas:}
    The \texttt{while} case relies on a previously mechanized
    \texttt{lemma\_2\_5}. This maps the textbook's informal reference to a
    formal transformation between the \texttt{if} branch and the
    \texttt{while} loop in the Natural Semantics (\texttt{ns}) inductive type.
    
    \item \textbf{Case Exhaustion (RQ1):}
    Lean's \texttt{match} and \texttt{cases} tactics ensure that the "inspecting
    the first step" strategy covers all possible SOS rules, mapping the
    textbook's "proceed by cases" directly to Lean's inductive elimination.
\end{itemize}

\subsubsection*{Completion of the Theorem}
\begin{paracol}{2}
\begin{leftcolumn}
    \begin{Lean_Code}
theorem sos_to_ns (S : Stmt) (s s' : State) :
  (@\(\langle\)@S, s@\(\rangle\)@ @\(\to_{sos}\)@* s') @\(\to\)@ @\(\mathbb{Z}\)@@\(\langle\)@S, s@\(\rangle\)@ @\(\to_{ns}\)@ s') := by
  intro h
  -- 1. Convert star to a k-step relation
  obtain @\(\langle\)@k, hk@\(\rangle\)@ := star_to_small_step_k h
  -- 2. Call a helper lemma that works via induction on k
  exact sos_k_to_ns S s s' k hk

theorem ns_equivalent_sos (S : Stmt) (s s' : State) :
  (@\(\langle\)@S, s@\(\rangle\)@ @\(\to_{ns}\)@ s') @\(\leftrightarrow\)@ (@\(\langle\)@ S, s@\(\rangle\)@ @\(\to_{sos}\)@* s') := by
  apply Iff.intro
  case mp => exact ns_to_sos S s s'
  case mpr => exact sos_to_ns S s s'
    \end{Lean_Code}
\end{leftcolumn}
\begin{rightcolumn}
\begin{quotation}
\noindent
Proof of Theorem 2.26:
For an arbitrary statement \(S\) and state \(s\), it follows from
Lemmas 2.27 and 2.28 that if \(\mathcal{S}_{ns}[\![S]\!]s = s'\) then \(\mathcal{S}_{sos}[\![S]\!]s = s'\) and vice versa.
This suffices for showing that the functions \(\mathcal{S}_{ns}[\![S]\!]\) and \(\mathcal{S}_{sos}[\![S]\!]\) must be equal:
if one is defined on a state s, then so is the other, and therefore if one is not
defined on a state s, then neither is the other.
\end{quotation}
\end{rightcolumn}
\end{paracol}
\vspace{1em}

\subsubsection*{Comparison and Result (Complete Proof)}
The proof of Theorem 2.26 (\texttt{ns\_equivalent\_sos}) is concise in both the
textbook and mechanized versions, but reveals important differences in how
mechanized proofs handle logical equivalences:
\begin{itemize}
    \item \textbf{Biconditional Construction:}
    The textbook simply states $\mathcal{S}_{ns}[\![S]\!] = \mathcal{S}_{sos}[\![S]\!]$
    and argues both directions. Lean must explicitly construct a biconditional
    ($\leftrightarrow$) using \texttt{Iff.intro}, which takes the forward
    direction (\texttt{mp}) and backward direction (\texttt{mpr}). 
    
    \item \textbf{Bridging Multi-Step Relations (RQ2):}
    The textbook implicitly equates the execution sequence $\Rightarrow^*$
    with the existence of a $k$-step derivation. Lean makes this implicit
    "shortcut" explicit through the \texttt{obtain} tactic, which extracts
    the step-count $k$ and the property $hk$ before applying Lemma 2.28. 
    
    \item \textbf{Delegating to Sub-Lemmas:}
    The proof in Lean does not repeat inductive arguments; instead, it applies
    \texttt{ns\_to\_sos} and \texttt{sos\_to\_ns}. This mirrors the textbook's
    reference to Lemmas 2.27 and 2.28, making the logical hierarchy and
    dependencies explicit to the kernel. 
    
    \item \textbf{Relational Interpretation (RQ1):}
    Both Lemmas are formulated as implications on evaluation relations
    ($\to_{ns}$ and $\to_{sos}^*$) rather than semantic functions.
    This is a key structural difference: the textbook's informal statement
    about functions'' is given a relational interpretation in Lean, which
    is more amenable to formal proof in a logic that handles partiality
    via relations.
    
    \item \textbf{Handling Partiality (RQ2):}
    As the textbook notes, equality of functions implies that if one is
    defined, the other is too. In Lean, this is handled by the relational
    definition where "undefined" simply means no $s'$ exists such that
    $(S, s) \to s'$. This bypasses the need for the "pedagogical shortcut"
    of discussing partial functions directly, instead relying on the
    completeness of the relation.
\end{itemize}

